

{\color{blue}
This paper proposes a new mixed integer linear mathematical programming (MILP) model, dubbed 3D-SCPTOP, for procurement transport planning in manufacturing supply chains. The model effectively integrates procurement and transport decisions while accounting for three-dimensional full-truck loading constraints and managing returnable transport items (RTIs). Thus, we optimise the material to procure for an assembly manufacturing site, RTIs closed-loop management,  and the involved transport decisions, including the truck loading configuration. Procurement transport decisions are based on the demand of assemblers, the supplier's inventory levels, and the dimensions of the truck / RTIs. The optimisation model aims to minimise transportation and inventory costs. Because the loading subproblem is NP-hard and is replicated once per loading configuration, the compact model is intractable at industrial sizes. We therefore reformulate it by Dantzig--Wolfe decomposition, taking the complete loaded truck as the unit of decision: a column is a fully specified, geometrically realisable truck load, so every placement, height, crush-limit, payload and full-truckload constraint holds inside it by construction, and the per-stack crush cap is evaluated on the actual product mix rather than on a worst-case bound. Columns are priced on demand by a constraint programming oracle and integrality is reached by depth-first branch and price, yielding MATH-3D-SCPTOP. Every plan the method returns is re-verified truck by truck, independently of the solver that produced it. The approach was validated with a case study in an automotive seating first-tier supplier, where it improves on the company's spreadsheet-based procedure in both transport and inventory terms. A comparison against a benchmark model that represents truck capacity volumetrically shows why the loading geometry cannot be abstracted away: the volumetric model returns cheaper plans whose truck loads cannot physically be loaded.
}

\keywords{ Procurement, inventory, transport planning,  full-truckload, 3D container loading }
\section{Introduction}

Optimizing the use of resources in manufacturing supply chains is crucial for boosting economic growth, improving social well-being, and minimizing environmental impact. Integrating transport, inventory, and production processes is a fundamental characteristic of supply chains \autocite{Stadtler2015SupplyPlanning} because joint decisions will improve supply chain performance\autocite{Diaz-Madronero2017ADecisions}.

Procurement refers to the process of acquiring goods and services in a supply chain \autocite{Chopra2015SupplyOperation}. One of the main goals of procurement is to obtain raw materials, parts or finished goods for a factory or distribution centre from a supplier. As mentioned in \autocite{Fleischmann2015TransportDistribution}, the responsibility for transport planning usually lies with the supplier. However, sometimes the manufacturer makes transport decisions of its supplies in, for example, the automotive industry. Furthermore, an essential asset in the automotive supply chain is not considered in traditional material requirement planning (MRP): returnable containers or returnable transport items (RTI), which are used as tertiary packaging. RTIs enable the storage and transportation of RM and FG through the supply chain. Usually, the manufacturing plant supplies empty RTIs to its supplier, while the supplier delivers full RTIs with the FG or RM. Therefore, during the procurement process, transport and inventory decisions should be integrated, including RTI management. Here, we focus on optimising the operational decisions made during procurement transport planning.

According to \autocite{Diaz-Madronero2014AChain}, using algorithm-based methods for procurement transport planning is more effective than relying on manual approaches based on planners' intuition. Manual planning can be short-sighted and lead to less efficient outcomes. This becomes even more significant when optimizing procurement transport tasks that aim to reduce the number of trucks and inventory levels, especially under 3D loading constraints where the precise placement of containers within trucks must also be determined (cite ).
Furthermore, the capacity constraints used in the scientific literature on procurement transport planning are based mainly on freight volume or the total number of loaded product units  {\autocite{Hashemi-Amiri2023IntegratedApproach} and \autocite{Diaz-Madronero2014AChain}.} The current research seeks to address this gap by introducing a new mixed-integer linear programming (MILP) model. This model uniquely combines procurement transportation with considerations for 3D loading, inventory management, and RTI management. Hereafter, we refer to this model formulation as 3D supply chain procurement transport operational planning (3D-SCPTOP). Additionally,  a matheuristic procedure (MATH-3D-SCPTOP)  is researched.



{\color{blue}

In this paper, we present four main contributions: (i) formulating a novel MILP model, 3D-SCPTOP, that integrates procurement transport, RTI management and inventory decisions with stack-based full-truck loading constraints; (ii) reformulating it by Dantzig--Wolfe decomposition over geometrically realisable truck loads, and solving it by branch and price with a constraint programming pricing oracle, giving MATH-3D-SCPTOP; (iii) showing that a volumetric representation of truck capacity, as used in the procurement transport literature, yields plans that are cheaper on paper but cannot be loaded in practice; (iv) validating and testing the proposal with a case study at a first-tier supplier of the automotive sector.
}


Readers are referred to the typology by  \autocite{Wascher2007AnProblems} for a global view of container loading problems, where the different dimensions of the loading problem are described. Furthermore, \autocite{Silva2019ExactProblems} summarise the exact methods used in the literature to solve container loading problems, and \autocite{Zhao2016INTERNATIONALAlgorithms} explore other solution approaches in the literature. Readers are also referred to the literature reviews of \autocite{Mosca2019IntegratedReview} and \autocite{Diaz-Madronero2015ADecisions}, which focus more on integrating inventory decisions and transport. Here, we jointly address container loading, procurement transport and inventory decisions. {\color{blue} Furthermore, we refer readers () for insights about the packaging methods used in supply chains.}

{\color{blue}
The paper is organised as follows. Section 2 provides a literature review related to the problem under study. Section 3 presents the problem description. Section 4 describes the 3D-SCPTOP MILP formulation. Section 5 develops the branch and price solution approach, MATH-3D-SCPTOP, together with its pricing subproblem and bounding arguments. Section 6 reports the computational experiments and Section 7 the case study. Finally, Section 8 concludes and puts forward further research lines.
}

\section{Literature Review}

The goal of this research is to integrate procurement transport operational planning decisions with 3D loading restrictions and RTI packaging. Based on the authors’ research, no previous work has been identified that meets this criterion. Therefore, three independent, but sometimes overlapping, knowledge bodies have been reviewed: procurement transport and inventory planning, container loading problems and RTI literature.

\subsection{Procurement transport and inventory planning}

Models that integrate transport and inventory costs during the procurement process are widely studied in the literature. As noted in \autocite{Langley2020SupplyPerspective}, transportation facilitates the movement of raw materials and finished products between suppliers,  warehouses, manufacturers, and retailers. Therefore, transport processes play a crucial role in supply chains. \autocite{Fleischmann2015TransportDistribution} distinguish between two main types of transportation within a supply chain: "procurement transport," which involves moving materials from suppliers to manufacturing sites, and the "distribution of goods," which concerns delivering finished products from factories to customers. This article focuses specifically on addressing the challenges related to procurement transportation.

There are two main dimensions of procurement transport problems according to the supply chain configuration and the applied transport method. Regarding supply chain configurations \autocite{Mosca2019IntegratedReview}: (i) single source to a single destination (SS); (ii) multiple sources to a single destination (MS);
(iii) single source to multiple destinations (SM); (iv) multiple sources to multiple destinations (MM). Source components determine the number of supplier locations that are taken into account, while destination components denote the number of supplied manufacturers. Additionally, different transport methods are identified in \autocite{Diaz-Madronero2014AChain}:milk run, less-than-truckload (LTL),  full-truckload (FTL) . We also analyze the formulation of the transportation capacity constraints in the scientific literature.

\autocite{Berman2006InboundCost} address the problem of optimising distribution strategies to deliver products from suppliers to plants, whose aim is to minimise transport, pipeline and plant inventory costs by using cross docks. A mixed integer non-linear programming (MINLP) formulation, a Lagrangian  relaxation heuristic,  a greedy heuristic, and a branch-and-bound algorithm are proposed to model and solve the problem.

\autocite{Sancak2011Multi-itemPolicy} present a MILP model to optimise procurement transport decisions in a SS supply chain configuration using LTL and FTL transport methods. The objective is to minimise transport and inventory holding costs, while fulfilling production plan requirements on a multiperiod planning horizon. The study explores the possibility of delaying LTL shipments to the subsequent time period, and utilising safety stocks whenever necessary, in a stochastic environment by employing data from a bus manufacturer. The findings demonstrate that delaying shipments leads to reduced holding and transport costs with a minimal stockout risk.


\autocite{Schoneberg2010AnNetworks} examine the concept of delivery profiles, which define the delivery frequency at a fixed rate, offering structured planning, minimizing management effort, and improving coordination in logistics. In this context, FTL transport is used when the goods' volume exceeds a vehicle's capacity, making direct shipment possible. LTL is applied when the volume is insufficient to fill a truck, and the goods are first sent to a consolidation center before reaching the final destination. In \autocite{Schoneberg2013ANetworks}, the proposed model accounts for demand uncertainty.

\autocite{Diaz-Madronero2014APlanning} present a  fuzzy  multi-objective  integer  linear programming (FMOILP) model that focuses on operational procurement transport planning in an automotive supply chain, which involves a single source and a single destination while utilising FTL. The developed model aims to minimise the number of trucks employed for procurement transport and the inventory level by considering the uncertainties related to transport capacity levels.

\autocite{Hein2016QuantitativeProblem} investigate the benefits of integrated procurement transport and production planning for supply and production processes by comparing a sequential approach, where production is firstly determined, followed by procurement transport decisions, plus an integrated approach in which decisions are jointly made. Numerical experiments reveal that the integrated planning approach leads to substantial cost savings.

\autocite{Diaz-Madronero2017ADecisions} present an novel MILP model that integrates  procurement, production and transport planning decisions within manufacturing systems. This model extends beyond traditional material requirement planning (MRP) systems by incorporating procurement transport decisions from multiple sources and a single destination into a unified optimization framework. It considers the capacities of both material and production resources, along with procurement transport planning decisions, by employing various shipping modes such as FTL, LTL, and milk-run.

\autocite{Chitsaz2019ARouting} address integrating production, inventory and procurement transport into a manufacturing context by defining this problem as an assembly routing problem (ARP). It introduces a novel MILP model for the ARP and proposes a three-phase decomposition matheuristic for its solution.

\autocite{Baller2022OptimizingApproach} examine the optimization of procurement transport in an automotive factory, focusing on first-tier suppliers at an aggregate level across multiple sources (570) to a single destination, using FTL, LTL, and courier express service (CES) as transportation methods. The proposed MILP approach includes a cyclic inventory constraint, truck utilization, and storage limitations.

\autocite{Hashemi-Amiri2023IntegratedApproach} propose an integrated approach to optimize a perishable food supply chain. The model combines supplier selection, procurement and distribution transport, production scheduling, and vehicle routing to improve supply chain efficiency and resilience. It handles demand uncertainties through a distributionally robust chance-constrained model. The model is reformulated as MILP by utilizing duality and linearization techniques.


\subsection{Procurement Transport with RTIs}
Some models further integrate the procurement and delivery of RTIs into the procurement transport literature. In this sense, RTIs are procured to the supplier, who loads the items and sends them to customer. We extend the SC definitions previously defined to contemplate the RTIs management strategies.

model a food closed-loop supply chain network design problem incorporating returnable transport items (RTIs), product perishability, facility location and demand uncertainty. A two-stage non-linear chance-constrained programming model minimizes fixed facility, production, transport, inventory holding and RTI costs while ensuring a target service level. It’s linearized via sample average approximation and distribution-free methods to capture forward RTI-based shipments and reverse flows for reuse. To solve large instances, an improved Lagrangian relaxation algorithm decomposes facility-location and flow decisions, outperforming CPLEX on solution quality and speed.


\autocite{SOYSAL2023101504} develop a tactical-level, multi-period closed-loop Vendor-Managed Inventory (VMI) system for RTIs in an agro-food SC. The SM SC integrates forward flows of full RTIs from a single vendor to multiple customers and reverse flows of empties back to the vendor for cleaning . All routing and inventory decisions are made centrally by the vendor under full information sharing. Sustainable plans in terms of amounts of energy spent and empty crates collected are devised and evaluated.


\autocite{GUO2024104339} model two-tier Container Management Centre (CMC) network . Empty RTIs are dispatched from central CMC hubs to upstream suppliers, loaded with products, carried by truck to downstream customers, then returned empty to regional CMCs for cleaning and storage before being shipped back to central hubs. A neural network is used for demand forecasting. A MILP is proposed that maximises total profit by optimally determining transport costs and the penalties applied to late RTIs returns.


\autocite{Zhang17012024} formulate a perishable-food inventory-routing MILP with multiple RTI types in a SM supply chain, integrating procurement of products and pickup of both filled and empty RTIs. Decisions are made on a operational level, optimally determining per period the quantities to pack, deliver, pick up, purchase and lose, as well as the vehicle routes in order to maximise the total profit of the SC that is equal to the total quality-dependent revenue minus the total costs involved, including inventory holding costs, product cost, transportation costs, and RTI purchasing and loss costs. The model is solved exactly via CPLEX  solver and approximately via a tailored Kernel Search matheuristic.


\subsection{Container loading}

Container loading problems encompass a subset of geometric assignment problems known as cutting and packing ($C\&P$) problems. They revolve around the allocation of n-dimensional orthotope items in larger n-dimensional orthotope containers. Here the aim is to optimise an objective function while taking into account two key geometric constraints: complete containment of items inside containers and lack of item overlap. A loading pattern serves as a formal representation of a solution to such an assignment problem.
Dyckhoff introduced in \autocite{Dyckhoff1990AProblems} a typology that unifies $C\&P$ problems using four different criteria: dimensionality of space, the objective of assignment, containers types and box type. Consequently, \autocite{Wascher2007AnProblems} develop a new typology that has become the standard in the $C\&P$ literature. We   briefly  describe  the  typology  used  in  \autocite{Wascher2007AnProblems}  based  on  five different criteria: (i) dimensionality of space; (ii) kind of assignment (input minimisation of containers for all small items and loading patterns; output maximisation to pack a subset of small items with limited container availability;) (iii) assortment of small item by distinguishing if they are either a weakly (or strongly) heterogeneous assortment or identical small items ; (iv) assortment of large objects (if it is either several large objects or one large object); (v) shape of small items (circular or of another shape). Table 1 links the standard $C\&P$ nomenclature and the criteria of the typology according to \autocite{Wascher2007AnProblems}.





The first mathematical programming model in a scientific journal to describe the container-loading problem was developed in \autocite{Chen1995EUROPEANProblem}, where the authors introduced an MILP model to address the 3D  input minimisation problem by considering rotations and containers of varying dimensions. However, it does not specify the use of weakly or strongly heterogeneous boxes or containers. Additionally, real-world constraints are modelled like, for example, weight distribution.

\autocite{Junqueira2012Three-dimensionalConstraints} propose an alternative MILP formulation to \autocite{Chen1995EUROPEANProblem} using binary encoding for box position, which utilises normal patterns to reduce possible box coordinates by considering multiples of boxes dimensions and ensuring optimal solutions. Additionally, loading and cargo constraints are developed to ensure the static and dynamic stability of the cargo. Boxes must be supported by other box top faces or the container floor. The same authors integrate further loading restrictions in \autocite{Junqueira2012MIP-basedConstraints}, where they integrate multidrop constraints, which address the challenge of packing a
truck with boxes for different customers at various destinations. The goal is to plan the container’s loading to minimise the time spent unloading and reloading boxes at each destination.

\autocite{Zhu2012AProblem} devise a tree search-based method to solve SLOPP and SKP with 3D constraints. Their approach involves evaluating free spaces by employing a look-ahead tree search and selecting optimal placements based on the best free space-block pairs.

\autocite{Goncalves2013AProblems} develop a biased random-key genetic algorithm (BRKGA) designed to solve SBSBPP. The algorithm addresses geometric constraints and introduces a maximal-space representation to handle available free spaces in bins. A novel placement procedure is hybridised with the BRKGA, which optimises the order of box packing and the parameters utilised by the placement procedure to offer an effective solution approach to an NP hard problem.

\autocite{Bortfeldt2013ConstraintsReview} identify constraints for $C\&P$, including item-related constraints (orientation constraints,   load bearing, loading priorities), container-related constraints (weight distribution and weight limit),  and cargo-related constraints (complete-shipment constraints) and load-related constraints ( complexity constraints or stability constraints). The study also highlights the achievements and deficits in container loading research that pertain to these constraints.

\autocite{2016Paquay} propose an MILP model for the MBSBPP.  One noteworthy constraint is the assumption that containers are truncated parallelepiped instead of rectangular boxes. They also consider constraints, such as box rotation, container load bearing, weight distribution, stability and fragility of items.

\autocite{Mahvash2018} introduce a novel resolution procedure for the SBSBPP through the column generation technique, where the columns that are generated are the loading configurations of a single bin. For the pricing develop problem they develop a heuristic to find new columns with a negative reduced cost.  This approach not only produces better results than the previous articles on the literature, but also gives new lower bounds.




\autocite{Silva2019ExactProblems} study SSLOP models with 3D constraints. These authors compare the different approaches that exist in the literature for model packing problems based on two fundamental approaches: discretised formulation \autocite{Junqueira2013AnConstraints} and continuous formulation \autocite{Chen1995EUROPEANProblem} with subsequent modifications. They also indicate that discretised formulation\autocite{Junqueira2012MIP-basedConstraints} performs the best for instances with similar sized boxes. However, it encounters memory limitations for large instances. For other sets of instances, the adapted formulations of \autocite{Chen1995EUROPEANProblem} and \autocite{2016Paquay} outperform other adaptations. Further research topics suggest including the design of new formulations by considering the limited approaches in the literature.



\autocite{Kurpel2020TheProblems} introduce techniques and exact approaches to solve input minimisation and output maximisation problems for container-loading problems. Their methods consider practical constraints, such as orientation, separation of boxes and vertical stability. A set partitioning method is also employed in order to obtain upper bounds in a  container-loading input minimisation problems. The authors adapt and evaluate four discretisation methods found in the existing literature and incorporate mathematical formulations to address the aforementioned constraints.



\autocite{Alonso2020AConstraints} develop an algorithm by the greedy randomised adaptive search procedure (GRASP) to tackle a practical container-loading problem with time-sensitive constraints. Their approach addresses dynamic stability constraints while stacking products on pallets and loading them onto trucks during different time periods. Their approach provides high-quality solutions and surpasses an integer programming formulation in efficiency and effectiveness terms.

\autocite{Ali2022On-lineApproaches} conduct a comprehensive literature review by identifying constraints and solution methods for offline problems with full knowledge of the items to be packed and those tackling the online problems, where items are packed immediately when they arrive, and no prior knowledge of an item is known. The study also identifies and classifies solution approaches for offline packing into six categories: exact methods, constructive heuristics; local search and metaheuristic methods; approximation algorithms; machine learning and hyper-heuristic methods; tree search methods.
\autocite{Harrath2022} proposed a thee stage heuristic based on layer generation to find good quality solutions  the SBSPP problem for large instance up to 993 bins and 200000 items taking into account load balancing constraints.


\autocite{Tresca2022} develop a matheuristic to solve the SBSBPP that includes items with different IDs and different categories, to be delivered to different clients. The algorithm creates monoitem, monocategory, and mixed layers which then are used to compute the final loading configurations of the bins.

\autocite{Li2022} develope a hybrid adaptaive large neighborhood search algoritm to solve the multiple container loading problem, which is tries pack a given set of items with the minimum number of containers while maximazing the space utilization of them. Therfore, it is both a MBSBPP and a MSSCSP. Their approach takes into account some practical applications such as supsension constraints, rotation and weight limit.

\autocite{Chen2023}  present a new approach to the 3D-BPP, introducing the 3D-BSDPP which focuses on optimizing both the design of bin sizes and the packing process. A novel Hybrid Biogeography-Based Optimization (HBBO) algorithm is proposed, integrating differential evolution and a local search strategy for enhanced performance.



One noteworthy trend in $C\&P$ is integrating loading problems with routing. This is known as capacitated vehicle-routing problem with 3D or 2D loading constraints (3L-CVRP or 2L-CVRP). \autocite{Junqueira2012Three-dimensionalConstraints} present a MILP model with binary encodings for 3L-CVRP by including loading, cargo and multidrop constraints. The combination of these two NP-hard problems requires using heuristics and metaheuristics for real-world instances, as seen in \autocite{Dominguez2016UsingFleet}, \autocite{Alonso2022TheRevision} and \autocite{Wei2015AConstraints}. Another noteworthy trend is to use reinforcement learning techniques to solve $C\&P$, as in \autocite{Que2023SolvingLearning} and \autocite{Zhao2022LearningPacking}.


\subsection{Research Gaps \& Methodology}


Here, we discuss the main differences between the 3D-SCPTOP model proposal and those reviewed by specifically focusing on procurement transport. The key distinction of the 3D-SCPTOP model compared to the existing literature is the incorporation of 3D constraints and vertical stability constraints to model transport capacity with a view to providing more realistic solutions and reducing the likelihood of infeasible outcomes.   Table \ref{table:ComProTra}  presents a comparison of inventory modelling, transport modelling, and supply chain configuration between the 3D-SCPTOP model proposal and those currently in the literature.
Hence,  the 3D-SCPTOP approach is based on a SS supply chain configuration, and takes into account the FTL transportation mode.
\begin{table}[h!]
  \begin{adjustbox}{width=\textwidth}
    \begin{tabular}{|c|c|c|c|c|c|c|l|c|c|c|}
      \hline
      Article                                                                            &
      \begin{tabular}[c]{@{}l@{}} SC \\ Configuration\end{tabular}                       &
      \begin{tabular}[c]{@{}l@{}}Decision\\  level\end{tabular}                          &
      \begin{tabular}[c]{@{}l@{}}Transportation \\ Mode\end{tabular}                     &
      \begin{tabular}[c]{@{}l@{}}Transportation \\ Capacity \\   Constraint\end{tabular} &
      \begin{tabular}[c]{@{}l@{}}Transportation\\  Cost\end{tabular}                     &
      \begin{tabular}[c]{@{}l@{}}Inventory\\  Capacity \\ Restriction\end{tabular}       & RTIs           &
      \begin{tabular}[c]{@{}l@{}}Inventory\\  Costs\end{tabular}                         &
      \begin{tabular}[c]{@{}l@{}}Safety \\ Stocks\end{tabular}                           &
      \begin{tabular}[c]{@{}l@{}}Modelling \\ Framework\end{tabular}                                                                                    \\
      \autocite{GUO2024104339}                                                           & MM             & TO & FTL & UH & X & X & S|NF & X &  & MINLP \\\hline \hline
      \autocite{Berman2006InboundCost}                                                   &
      MM                                                                                 &
      T                                                                                  &
      FTL|LTL                                                                            &
      V                                                                                  &
      X                                                                                  &
                                                                                         &                &
      X                                                                                  &
                                                                                         &
      MINLP                                                                                                                                             \\ \hline
      \autocite{Sancak2011Multi-itemPolicy}                                              &
                                                                                         &
      O                                                                                  &
      LTL|FTL                                                                            &
      V                                                                                  &
      X                                                                                  &
                                                                                         &                &
      X                                                                                  &
      X                                                                                  &
      MILP                                                                                                                                              \\ \hline
      \autocite{Schoneberg2010AnNetworks}                                                &
      MS                                                                                 &
      T                                                                                  &
      FTL|LTL                                                                            &
      V                                                                                  &
      X                                                                                  &
      X                                                                                  &                &
      X                                                                                  &
                                                                                         &
      MILP                                                                                                                                              \\ \hline
      \autocite{Schoneberg2013ANetworks}                                                 &
      MS                                                                                 &
      O                                                                                  &
      FTL|LTL                                                                            &
      V                                                                                  &
      X                                                                                  &
      X                                                                                  &                &
      X                                                                                  &
                                                                                         &
      MILP                                                                                                                                              \\ \hline
      \autocite{Diaz-Madronero2014APlanning}
                                                                                         &
      SS                                                                                 &
      O                                                                                  &
      FTL                                                                                &
      UH                                                                                 &
      X                                                                                  &
                                                                                         &                &
      X                                                                                  &
      X                                                                                  &
      FMOILP                                                                                                                                            \\ \hline
      \autocite{Hein2016QuantitativeProblem}                                             &
      MS                                                                                 &
      O                                                                                  &
      MR                                                                                 &
      V                                                                                  &
      X                                                                                  &
                                                                                         &                &
      X                                                                                  &
                                                                                         &
      MILP                                                                                                                                              \\ \hline
      \autocite{Diaz-Madronero2017ADecisions}                                            &
      MS                                                                                 &
      T                                                                                  &
      FTL|LTL|MR                                                                         &
      V                                                                                  &
      X                                                                                  &
      X                                                                                  &                &
      X                                                                                  &
      X                                                                                  &
      MILP                                                                                                                                              \\ \hline
      \autocite{Meyer2018TransportIndustry}                                              & MS
                                                                                         & T|O
                                                                                         & MR|FTL
                                                                                         & V
                                                                                         & X
                                                                                         & X
                                                                                         &                & X
                                                                                         &
                                                                                         & MILP
      \\ \hline
      \autocite{Chitsaz2019ARouting}                                                     &
      MS                                                                                 &
      O|T                                                                                &
      MR                                                                                 &
      V                                                                                  &
      X                                                                                  &
      X                                                                                  &                &
      X                                                                                  &
                                                                                         &
      MILP                                                                                                                                              \\ \hline
      \autocite{Baller2022OptimizingApproach}                                            & MS
                                                                                         & O
                                                                                         & MR|FTL|LTL|CES
                                                                                         & V|W
                                                                                         & X
                                                                                         & X
                                                                                         &                &
                                                                                         & X
                                                                                         & MILP
      \\ \hline
      \autocite{Hashemi-Amiri2023IntegratedApproach}                                     &
      MS                                                                                 &
      T                                                                                  &
      MR                                                                                 &
      V                                                                                  &
      X                                                                                  &
                                                                                         &                &
                                                                                         &
                                                                                         &
      CCP                                                                                                                                               \\ \hline
      3D-SCPTOP                                                                          &
      SS                                                                                 &
      O                                                                                  &
      FTL                                                                                &
      3D-Constraints                                                                     &
      X                                                                                  &
      X                                                                                  &                &
      X                                                                                  &
      X                                                                                  &
      MILP                                                                                                                                              \\ \hline
      \autocite{Zhang17012024}                                                           & MS             & TO & MR  & UH & X & X & M|NF & X &  & ILP   \\\hline
    \end{tabular}
  \end{adjustbox}
  \begin{tablenotes}
    \scriptsize
    \item  \textbf{Abbreviations.}  \textbf{T} (tactical), \textbf{O} Operational; \textbf{FTL} (full-truck load), \textbf{LTL} (less-than-truck load), \textbf{MR} (milk-run), \textbf{CES} (courier express service); \textbf{V} (volumetric), \textbf{W} (weight), \textbf{UH} Unit Holding;   \textbf{CCP} chance-constrained programming
  \end{tablenotes}
  \caption{Comparison of procurement transport models}
  \label{table:ComProTra}
\end{table}

On the other hand, the 3D-SCTOP model diverges from those proposed in the packaging literature in several aspects. Firstly, the loading problem takes into account different time periods. Therefore, the specific container (small items or boxes in the $C\&P$ literature) that we utilise are yet to be defined, as the choice of container depends on the inventory levels of each product.
Also, the objective function depends on inventory and transport costs, which are inherently connected to the loading problem at hand. Thefore, we simultaneously solve an input minimisation  and output maximisation problem. Both of these are a result of transport minimisation and inventory costs. Input minimisation comes because we want to find the minimum number of trucks needed. Meanwhile, output maximization comes from the FTL and the transport costs.

\begin{table}
  \begin{adjustbox}{width=\textwidth}
    \begin{tabular}{|c|c|c|c|c|c|c|c|} \hline
      Article                           & Problem type     & Resolution procedure & Orientation & Weight & Load Balance & Stability & Time-Depenent \\
      \hline
      \autocite{Zhu2012AProblem}        & SLOP|SKP         & TSM                  & X           & X      &              &           &               \\ \hline
      \autocite{Goncalves2013AProblems} & SBSBPP           & BRKGA                & X           & X      &              & X         &               \\ \hline
      \autocite{2016Paquay}             & SBSBPP           & EM                   & X           & X      & X            & X         &               \\ \hline
      \autocite{Mahvash2018}            & SBSBPP           & CG                   & X           & X      &              &           &               \\ \hline \autocite{Alonso2020AConstraints}& SSSCSP& GRASP& X& X& X& X&X\\ \hline
      \autocite{Tresca2022}             & SBSBPP           & MH                   & X           & X      & X            & X         &               \\ \hline
      \autocite{Harrath2022}
                                        & SBSBPP           & H                    & X           & X      & X            &           &               \\ \hline
      \autocite{Li2022}                 & SBSBPP|MSSCSP    & HALNs                & X           & X      &              & X         &               \\
      \hline
      \autocite{Chen2023}               & BSDPP            & HBBO                 & X           & X      &              &           &               \\
      \hline
      3D-SCPTOP                         & SSSCSP \& MILOPP & MH                   & X           & X      & X            & X         & X             \\ \hline
    \end{tabular}
  \end{adjustbox}
  \begin{tablenotes}
    \scriptsize
    \item  \textbf{Abbreviations.}  \textbf{T}SM (Tree search method), \textbf{BRKGA} Biased Randomized Jey Genetic Algorithm; \textbf{EM} (Exact Metods), \textbf{MH} (MatHeuristic), \textbf{CG} (Column Generation), \textbf{H} (Heuristic); \textbf{HCLP} (Hybrid Adaptive Large Neighborhood Search), \textbf{HBBO} (Hybrid Biogeography-Based Optimization)
  \end{tablenotes}
  \label{tab:my_label}
  \caption{Comparison of C\&P problems}
\end{table}

% \begin{figure}
%     \centering
%     \includegraphics[scale=0.5]{FiguraPaperv2-cropped.pdf}
%     \caption{Caption}
%     \label{fig:enter-label}
% \end{figure}


Based on the literature review and the identified research gaps, our proposal is structured according to the steps depicted in Figure \ref{fig:tikz1}.

\begin{figure}[h!]
  \centering
  \resizebox{0.85\columnwidth}{!}{
    \begin{tikzpicture}[node distance=2cm, auto, thick]

      % Define styles
      \tikzstyle{block} = [rectangle, draw, fill=blue!40, text width=14em, text centered,  minimum height=3em]
      \tikzstyle{sideblock} = [rectangle, draw, fill=white, text width=10em, text centered, rounded corners, minimum height=2em]
      \tikzstyle{grayblock} = [rectangle, draw, fill=gray!20, text width=13em, text centered, rounded corners, minimum height=2.5em]
      \tikzstyle{line} = [draw, -Latex]
      \tikzstyle{plainline} = [draw]
      \tikzstyle{dashedline} = [draw, dashed, gray, thick, -Latex]





      % Central nodes
      \node [block] (literature) {Literature Review};
      \node [grayblock, below=0.35cm of literature] (research) {Research Gap};
      \node [block, below=0.35cm of research] (problem) {Problem Statement};
      \node [grayblock, below=0.35cm of problem] (3dMilp) {3D-SCPTOP};
      \node [grayblock, below=0.35cm of 3dMilp] (3dMH) {MH-3D-SCPTOP};
      \node [block, below=0.35cm of 3dMH] (experiments) {Computational Experiments};
      \node [grayblock, below=0.35cm of experiments] (procedure) {Resolution \\ Procedure Comparison};
      \node [grayblock, below=0.35cm of procedure] (comparison) {Model Comparison};
      \node[ aux_gene_group,fit=(comparison) (procedure)] (gr1) {};

      \node [block, below=0.35cm of comparison] (case) {Case Study \& \\ Managerial Implications};

      % Side nodes
      \node [sideblock, above=1cm of literature] (model) {Modelling \& \\ Resolution Appoach};
      \node [sideblock, left=1cm of literature] (procurement) {Procurement Transport};
      \node [sideblock, right=1cm of literature] (3D) {3D Container Loading};

      % Draw edges for side nodes (with dashed gray lines and arrows)
      \path [dashedline] (3D) -- (literature);
      \path [dashedline] (model) -- (literature);
      \path [dashedline] (procurement) -- (literature);

      % Side nodes
      \node [sideblock, left=1cm of 3dMH] (MILP) {MILP};
      \node [sideblock, right=1cm of 3dMH] (MetaHeuristic) {Metaheuristic};

      % Draw edges for side nodes (with dashed gray lines and arrows)
      \path [dashedline] (MILP) -- (3dMH);
      \path [dashedline] (MetaHeuristic) -- (3dMH);

      % % Side nodes
      % \node [sideblock, above left=1.25cm and 1cm of comparison] (small) {Small Instances};
      % \node [sideblock, below=0.5cm of small] (medium) {Medium Instances};
      % \node [sideblock, below=0.5cm of medium] (large) {Large Instances};

      % Side nodes
      \node [sideblock, left=0.5cm of gr1] (medium) {Medium Instances};

      \node [sideblock, above=0.5cm of medium] (small) {Small Instances};
      \node [sideblock, below=0.5cm of medium] (large) {Large Instances};

      \coordinate (center13) at ($(procedure.west)!1/3!(procedure.south west)$);
      \coordinate (center12) at ($(procedure.west)$);
      \coordinate (center23) at ($(procedure.west)!2/3!(procedure.north west)$);


      % Draw edges for side nodes (with dashed gray lines and arrows)
      \path [dashedline] (small.east) -- (center23) ;
      \path [dashedline] (medium.east)  -- (center12) ;
      \path [dashedline](large.east)  -- (center13) ;


      \coordinate (ccenter13) at ($(comparison.west)!1/3!(comparison.south west)$);
      \coordinate (ccenter12) at ($(comparison.west)$);
      \coordinate (ccenter23) at ($(comparison.west)!2/3!(comparison.north west)$);


      % Draw edges for side nodes (with dashed gray lines and arrows)
      \path [dashedline] (small.east) -- (ccenter23) ;
      \path [dashedline] (medium.east)  -- (ccenter12) ;
      \path [dashedline](large.east)  -- (ccenter13) ;




      % Draw edges for central nodes and gray nodes
      \path [plainline] (literature)  -- (research);
      \path [line] (research) -- (problem);
      \path [plainline] (problem)  -- (3dMilp);
      \path [plainline] (3dMilp)  -- (3dMH);
      \path [line] (3dMH) -- (experiments);
      \path [plainline] (experiments)  -- (procedure);
      \path [plainline] (procedure)  -- (comparison);
      \path [line] (comparison) -- (case);






      % Side nodes
      \node [sideblock, right=1cm of comparison] (reference) {With respect \autocite{Diaz-Madronero2014APlanning}};

      % Draw edges for side nodes (with dashed gray lines and arrows)
      \path [dashedline] (reference) -- (comparison);







    \end{tikzpicture}
  }
  \caption{Research Methodology}
  \label{fig:tikz1}
\end{figure}

\FloatBarrier




{
% =====================================================================
% 3D-SCPTOP -- stack-based reformulation (cleaned)
% Drop-in replacement for Sections 3 and 4 of the original manuscript,
% plus the C&P comparison row of Table 2 and noted prose edits.
%
% Index conventions used below:
%   g in G = {s, p}         node (supplier s, plant p)
%   n in N                  product
%   r in R                  RTI type
%   h in H                  stack type  (renamed from the earlier draft
%                           to avoid clashing with supplier node s)
%   i, j in I               stack-slot indices within an LC
%   l in L                  loading configuration
%   k in K                  truck
%   t in T                  planning period
%
% Runtime variables whose direction is fixed by definition drop the
% direction subscript:
%   V_{klnt}            -- products on plant-bound truck (implicit g=p)
%   V^{slot}_{kilnt}    -- products in slot i of plant-bound truck
%   V^{e}_{klrt}        -- folded RTIs on supplier-bound truck
% =====================================================================


% =====================================================================
%               S E C T I O N   3   (R E V I S E D)
% =====================================================================

\section{Problem Description}

The objective is to minimise the transportation and inventory holding
costs of a manufacturing site by integrating procurement transport,
RTI management and inventory decisions with full-truckload (FTL)
loading constraints expressed in a stack-based form. The supply chain
under study has a single source and a single destination. Products are
shipped in different RTI types. RTIs are not placed individually in
three dimensions but are grouped into vertical \emph{stacks} whose
2D footprints are placed on the truck floor. Each stack type has a
fixed footprint and is compatible with a subset of RTI types; each
layer of a stack is homogeneous (all RTIs in a layer share a single
type), while different layers of the same stack may carry different
compatible RTI types. The total stack height is bounded by the truck
height and the weight borne by the bottom of each stack is bounded by
a crush-limit cap. Trucks carry full RTIs from the supplier to the
plant and return folded empties from the plant to the supplier;
folding affects RTI height, and hence the number of layers that fit in
a stack, but not the stack footprint.

The assembler's demand drives the procurement process because it
determines the optimal quantity of each product to procure from the
supplier to the manufacturing site. Raw materials and assembled parts
are produced at the supplier, packed in their corresponding RTI, and
transported under FTL. Each RTI carries a single product type at any
given time, and each product is shipped in a single RTI type. Once
full RTIs arrive at the plant, the products are consumed and the RTIs
are emptied, folded and returned to the supplier.

{\color{blue}
Stacks must not overlap and must be contained within the truck's
floor; the truck's total weight is bounded; and the per-stack weight
cap is respected with \emph{period-exact} knowledge of the products
loaded in each slot, so that crush-limit tests reflect the actual
period's load rather than a worst-case bound. Stacks may be rotated
by~90$^\circ$ on the truck floor. Because stacks are self-supporting
by construction, no additional vertical-stability constraints are
required. As FTL is the mode of transport, a dispatched truck must be full ---
but a truck may be full either because it has run out of space or
because it has run out of payload, so the requirement is imposed on
whichever of the two resources binds.
Minimum coverage stock is enforced on future demand.
}

{\color{blue}
\paragraph{Initial data.}
\begin{itemize}
    \item \textbf{Product data:} initial inventory, periodic demand
    (in full RTIs) at the plant and at the supplier, product-to-RTI
    mapping, weight of each full RTI.
    \item \textbf{RTI data:} height when full and when folded, weight
    of a folded RTI, volume.
    \item \textbf{Stack data:} footprint $(p_h, q_h)$ of each stack
    type $h\in H$, set of compatible RTI types $R(h)$, and, for each
    $r\in R(h)$, number of RTIs per layer $n_{hr}$.
    \item \textbf{Transport data:} number of available trucks, truck
    length/width/height, truck weight capacity, per-stack weight cap,
    minimum truck utilisation on the binding resource.
\end{itemize}
}

\paragraph{To determine.}
\begin{itemize}
    \item A set of loading configurations (LCs), each specifying the
    stack-type assignment, 2D placement and orientation of its slots,
    and the number of layers of each compatible RTI type in each slot.
    \item The number of trucks used per day and the LC they employ.
    \item The number of full RTIs of each product loaded in each slot
    per period (plant-bound trucks) and the number of folded RTIs of
    each type loaded per truck (supplier-bound trucks).
    \item The resulting inventory of full and folded RTIs at both
    nodes.
\end{itemize}

{\color{blue}
\paragraph{Subject to.}
\begin{itemize}
    \item 2D non-overlap and containment of stacks on the truck floor,
    with allowed 90$^\circ$ rotation.
    \item Per-stack height budget.
    \item Per-stack crush-limit weight cap, period-exact on plant-bound
    trucks and deterministic on supplier-bound trucks.
    \item Per-truck weight capacity.
    \item Minimum truck utilisation, on volume or on payload,
    whichever binds.
    \item Minimum coverage stock.
\end{itemize}
}

\paragraph{In order to:}
\begin{itemize}
    \item Minimise total transportation cost.
    \item Minimise total inventory holding cost.
\end{itemize}


% =====================================================================
%               S E C T I O N   4   (R E V I S E D)
% =====================================================================

\input{sec_model_blue}

\begingroup\color{blue}
% =====================================================================
%  Branch and Price Solution Approach for 3D-SCPTOP
% =====================================================================
\section{Branch and Price Solution Approach}
\label{sec:math}

There are three structural difficulties with the 3D-SCPTOP model of
Section \ref{sec:model}. The first is that a priori we do not know how
many loading configurations are needed, so $|L|$ must be
overestimated; every configuration that turns out to be unnecessary
still carries its full block of placement variables. The second is
that the model replicates a complete two-dimensional packing problem
$|G|\cdot|L|$ times, and each replica contributes
$4\binom{|I|}{2}$ big-$M$ disjunctions
\eqref{sb:eq:nonov-x1}--\eqref{sb:eq:nonov-y2} whose linear
relaxation is weak. The third is that the configurations are mutually
interchangeable, so the search tree is massively symmetric.

It would overstate the case to call the resulting model intractable in
general, and Section \ref{sec:results} shows why: with a good
incumbent supplied, instances spanning five periods are solved to
proven optimality in seconds, even at sixty thousand variables. What
defeats it is the planning horizon. The placement block is replicated
per configuration rather than per period, but the truck-assignment and
inventory blocks grow with $|T|$, and it is their interaction with the
symmetric configuration design that explodes: at ten periods the model
returns no feasible solution within half an hour. Operational
procurement planning lives at horizons of weeks, so that is the regime
the method has to serve, and it is the regime in which a decomposition
earns its place. A column generation scheme based on branch and price
is therefore developed.

The key components of the branch and price algorithm are the master
problem (MP), the restricted master problem (RMP), the pricing
problem (PP) and the branching policy. The MP is a reformulation of
the model that exploits its decomposable structure. Inspecting
Section \ref{sec:model}, every geometric constraint
\eqref{sb:eq:slot-active}--\eqref{sb:eq:util} carries a single
configuration index $l$, and every loading constraint
\eqref{sb:eq:slot-cap}--\eqref{sb:eq:stack-w-sup} carries a single
truck index $k$ and a single period $t$; the transport and inventory
block sees them \emph{only} through the aggregate flows $Q_{nt}$ and
$Q_{rt}$ of \eqref{sb:eq:flow-n}--\eqref{sb:eq:flow-r}. The model is
therefore block-angular, with one independent block per dispatched
truck. This suggests taking as the unit of decision neither the
individual RTI nor the loading configuration, but the
\emph{complete loaded truck}.

The MP is typically too large to be solved directly because it
involves a huge number of columns representing all possible truck
loads. To manage this complexity we consider only a smaller subset of
columns, which leads to the RMP. Solving the RMP provides an optimum
for the included subset, but not necessarily for the entire column
set. RMP integrality constraints are therefore relaxed to obtain dual
variables, and the PP uses those duals to identify columns of
negative reduced cost that should be added to the RMP. If no such
column exists, the integrality-relaxed MP has been solved to
optimality. If all the variables associated with the solution are
integers, an integer solution to the current branch has been found.
If not, the branch must be divided by adding further constraints.
Figure \ref{fig:B&P3D} exemplifies the algorithm.

\begin{figure}[h!]
    \centering
    \resizebox{0.8\columnwidth}{!}{%
\begin{tikzpicture}[node distance=2cm, auto, thick]
    \tikzstyle{block} = [rectangle, draw=white, fill=white, minimum height=2em,
                         text width=17em, text centered]
    \tikzstyle{line} = [draw, -Latex]

    \node [block] (MP) {\textbf{Master Problem}\\[2pt]
        \footnotesize all feasible truck loads $P^{p}\cup P^{s}$};
    \node[block, below=0.9cm of MP] (RMP) {\textbf{Restricted Master Problem}\\[2pt]
        \footnotesize loads generated so far, $\lambda\ge 0$};
    \node[block, below=0.9cm of RMP] (PP) {\textbf{Pricing Problem}\\[2pt]
        \footnotesize design one truck: 2D placement,\\
        \footnotesize layering and product assignment};

  \draw[bend right,left,->,align=left] (MP.west) to
        node {A reduced number \\ of columns} (RMP.west);
  \draw[bend right,left,->,align=left] (RMP.west) to
        node {Dual prices $\pi_{ct}$, $\nu_{gt}$} (PP.west);
  \draw[bend left,right,<-,align=left] (RMP.east) to
        node {If reduced cost is negative \\ add the truck load as a column; \\
              if not, branch} (PP.east);
  \draw[bend right,below,->,align=center] (PP.west) to
        node {Maximise $\sum_{c}\pi_{ct}a_{c}$ \\ subject to
              \eqref{eq:PPtype}--\eqref{eq:PPsymm}} (PP.east);
  \draw[bend left,below,<-,align=center] (RMP.west) to node {} (RMP.east);
\end{tikzpicture}
}
\caption{Branch and price scheme for 3D-SCPTOP. Every geometric,
height, crush-limit, payload and full-truckload constraint is enforced
inside the pricing problem, so each column is a physically loadable
truck.}
\label{fig:B&P3D}
\end{figure}


\subsection{Formulation of the Master Problem and the Restricted
Master Problem}
\label{sec:mp}

The key idea of the set-partitioning formulation for 3D-SCPTOP is to
enumerate all the feasible \emph{truck loads} with their associated
costs. A plant-bound truck load is a vector
$\mathbf{a}=(a_{n})_{n\in N}$ of full RTIs for which there exist a
slot count $m\le|I|$, stack types $h_{1},\dots,h_{m}\in H$,
orientations, bottom-left coordinates $(x_{i},y_{i})$, layer counts
$\lambda_{ir}$ and per-slot product loads $v_{in}$ such that
\begin{align}
&\text{the $m$ footprints, rotated as chosen, do not overlap and lie
   inside } pK\times qK,
   \label{eq:col-geom}\\
&\textstyle\sum_{r\in R(h_{i})} r_{pr}\lambda_{ir}\;\le\;rK
   &&\forall i\le m,
   \label{eq:col-height}\\
&\textstyle\sum_{n\in M(r)} v_{in}\;=\;n_{h_{i}r}\lambda_{ir}
   &&\forall i\le m,\;r\in R(h_{i}),
   \label{eq:col-assign}\\
&\textstyle\sum_{n\in N} w_{n}v_{in}\;\le\;W^{\max}
   &&\forall i\le m,
   \label{eq:col-crush}\\
&\textstyle\sum_{n\in N} w_{n}a_{n}\;\le\;pK^{w},
   \qquad
   \max\Bigl\{\tfrac{\sum_{n}\gamma_{p\rho(n)}a_{n}}{\Gamma_{K}},\;
              \tfrac{\sum_{n}w_{n}a_{n}}{pK^{w}}\Bigr\}\;\ge\;\alpha,
   \label{eq:col-truck}\\
&a_{n}=\textstyle\sum_{i\le m} v_{in}
   &&\forall n\in N.
   \label{eq:col-agg}
\end{align}
Let $P^{p}$ be the set of all such loads, and $P^{s}$ the analogous
set of supplier-bound loads $\mathbf{b}=(b_{r})_{r\in R}$, built with
folded heights $r_{sr}$, folded weights $w^{\mathrm{f}}_{r}$ and
folded volumes $\gamma_{sr}$; there no product assignment is needed,
because the weight of a folded RTI depends only on its type.

Two features of this definition matter. First, a column carries the
\emph{per-slot} product assignment $v_{in}$, so the crush limit
\eqref{eq:col-crush} is evaluated on the product mix actually riding
in that slot. The column therefore reproduces the period-exact cap
\eqref{sb:eq:stack-w-plant} of the compact model, and not the
conservative surrogate $\bar w_{r}=\max_{n\in M(r)}w_{n}$ that a pool
of capacity-only configurations is forced to use. Second, a column is
a feasibility certificate: the placement
$(h_{i},x_{i},y_{i},\text{rotation})$ witnessing
$\mathbf{a}\in P^{p}$ is retained, so any solution can be handed to a
loading crew as a drawing and re-verified independently of the solver
that produced it.

Constraint \eqref{eq:col-truck} is the same disjunction as
\eqref{sb:eq:util}--\eqref{sb:eq:utilw}, written on the load a truck
actually carries rather than on the capacity its configuration
provides. Section \ref{sec:results} shows that both regimes occur
across the benchmark: at $v=60$ the trailer fills by volume, at
$v=120$ it fills by weight.

Let $\lambda^{g}_{pt}\in\mathbb{Z}_{\ge0}$ count the trucks
dispatched in direction $g$ and period $t$ carrying load $p$, and let
$a^{p}_{n}$ and $a^{p}_{r}$ be the RTIs of product $n$, respectively
of type $r$, that load $p$ contains. Every truck costs $cK$. The MP
of our problem is
\begin{align}
    \min\;\; cK \sum_{g\in G}\sum_{t\in T}\sum_{p\in P^{g}}\lambda^{g}_{pt}
    \;+\;\sum_{g\in G}\sum_{n\in N}\sum_{t\in T} ic_{gn} I_{gnt}
    \label{eq:MPObjective}
\end{align}
subject to
\begin{align}
    \sum_{p\in P^{p}} a^{p}_{n}\lambda^{p}_{pt} = Q_{nt}
      \quad & \forall n\in N,\;t\in T \label{eq:MPlinkp}\\
    \sum_{p\in P^{s}} a^{p}_{r}\lambda^{s}_{pt} = Q_{rt}
      \quad & \forall r\in R,\;t\in T \label{eq:MPlinks}\\
    \sum_{p\in P^{g}} \lambda^{g}_{pt} \le |K|
      \quad & \forall g\in G,\;t\in T \label{eq:MPfleet}\\
    \text{\eqref{sb:eq:inv-plant-main}--\eqref{sb:eq:safety}}
      \quad & \text{(inventory balances and coverage stock)}
      \nonumber\\
    \lambda^{g}_{pt}\in\mathbb{Z}_{\ge0}
      \quad & \forall g\in G,\;p\in P^{g},\;t\in T. \nonumber
\end{align}
Objective \eqref{eq:MPObjective} minimises transport plus inventory
holding cost. Constraints
\eqref{eq:MPlinkp}--\eqref{eq:MPlinks} state that the flow of full
and folded RTIs is exactly what the dispatched trucks carry, and
\eqref{eq:MPfleet} caps the fleet available per direction and period.

MP and 3D-SCPTOP have the same optimal value whenever $|L|$ is not
restricted a priori. Given a 3D-SCPTOP solution, each used truck
$(k,t)$ induces a load vector satisfying
\eqref{eq:col-geom}--\eqref{eq:col-agg}, obtained by reading the slot
data off its configuration and setting
$v_{in}=V^{\mathrm{slot}}_{kilnt}$; counting trucks by load vector
gives a MP solution of equal cost. Conversely, introducing one
configuration per distinct load used by a MP solution and assigning
trucks accordingly reconstructs a 3D-SCPTOP solution of the same
cost. This equivalence is what distinguishes the present approach
from a pool-based method such as PLC-SCPTOP: a predefined pool is a
\emph{restriction} and can only yield upper bounds of unknown
quality, whereas MP is the same problem written differently, so any
bound derived from it is a bound on the original.

The RMP replaces $P^{g}$ by a subset $\bar P^{g}$ and relaxes
integrality,
\begin{align}
    \min\;\; cK \sum_{g\in G}\sum_{t\in T}\sum_{p\in \bar P^{g}}\lambda^{g}_{pt}
    \;+\;\sum_{g\in G}\sum_{n\in N}\sum_{t\in T} ic_{gn} I_{gnt}
    \label{eq:RMPObjective}
\end{align}
subject to \eqref{eq:MPlinkp}--\eqref{eq:MPfleet} restricted to
$\bar P^{g}$, with $\lambda^{g}_{pt}\ge0$. Dual values are obtained
by solving the RMP. Let $\pi_{nt}$ and $\pi_{rt}$ be the duals of
\eqref{eq:MPlinkp}--\eqref{eq:MPlinks} and $\nu_{gt}\le0$ those of
\eqref{eq:MPfleet}. The reduced cost associated with a plant-bound
load $p\in P^{p}$ dispatched in period $t$ is
\begin{equation}
\bar c(p,t)\;=\;cK-\sum_{n\in N}\pi_{nt}a^{p}_{n}-\nu_{pt},
\label{eq:redcost}
\end{equation}
and analogously in the supplier direction. A load with negative
reduced cost can improve the current RMP solution. For each iteration
we solve the subproblem of finding the load of minimum reduced cost;
loads with negative reduced cost are added as new columns, the RMP is
re-solved to generate new duals, and the process repeats until no
further improving column is discovered.

\subsection{Lower bounds}
\label{sec:bounds}

Care is needed here, because the natural candidate is not a bound at
all. The RMP is a \emph{restriction} of the MP: it optimises over
$\bar P^{g}\subseteq P^{g}$, so
$z_{\mathrm{RMP}}\ge z_{\mathrm{MP}}$. Reporting the RMP value as a
lower bound is therefore valid only once column generation has
converged, and wrong before then. We use two bounds accordingly.

\paragraph{Certified bound.}
When the PP has \emph{proven} that
$\zeta_{gt}:=\max_{p\in P^{g}}\sum_{c}\pi_{ct}a^{p}_{c}\le cK-\nu_{gt}$
for every $g$ and $t$, no column outside $\bar P^{g}$ has negative
reduced cost, so the current duals are feasible for the dual of the
full MP. Weak duality then gives
$z_{\mathrm{RMP}}=z_{\mathrm{MP}}\le z^{\star}_{\text{3D-SCPTOP}}$.
We report such bounds as \emph{certified}; the gap against the final
integer solution is then a genuine optimality gap, and not a gap with
respect to a pool.

\paragraph{Lagrangian bound.}
Certification requires the two-dimensional packing subproblem to be
closed to proven optimality, which is not always possible inside a
sensible time limit. Dualising
\eqref{eq:MPlinkp}--\eqref{eq:MPlinks} with the current $\pi$ leaves a
subproblem in which each column enters with reduced cost
\eqref{eq:redcost}; since the number of trucks an optimal solution can
dispatch in total is bounded by $\bar X$, the correction
\begin{equation}
z_{\mathrm{MP}}\;\ge\;z_{\mathrm{RMP}}
   \;+\;\bar X\cdot\min_{g\in G,\,t\in T}\;
   \min\bigl\{0,\;cK-\zeta_{gt}-\nu_{gt}\bigr\}
\label{eq:lagbound}
\end{equation}
holds at every iteration, converged or not. Two choices make
\eqref{eq:lagbound} usable rather than vacuous. The value substituted
for $\zeta_{gt}$ is the solver's own dual bound on the PP, never its
incumbent, so a subproblem stopped early still yields a valid
verdict. And $\bar X$ must be a genuine bound on the fleet actually
used: the textbook form charges $|K|$ separately for every direction
and period, which here is roughly an order of magnitude too
pessimistic and drove the reported gaps above $80\%$. Because
transport cost alone cannot exceed the objective, any incumbent
$\bar z$ gives $\bar X=\bar z/cK$, which is both valid and tight
enough to be informative.

\paragraph{What we report.}
We report the certified bound, and when certification fails we report
no bound at all. The Lagrangian correction \eqref{eq:lagbound} is
computed and recorded, but not published. In our implementation it
repeatedly returned values above solutions we could exhibit: on one
instance it gave $21552.23$ while the pool of that same run admits a
slack-free plan, verified truck by truck, costing $21502.80$. What
makes this dangerous rather than merely loose is that no check inside
the run detects it --- the run's own restricted master lies above the
true optimum too, so the estimate looks consistent from the inside,
and only an externally exhibited plan exposes it. A quantity that
cannot be falsified from within the method is not one we are willing
to print as a bound.

\paragraph{Two ways a certificate can be void.}
A certificate is only as good as the exhaustiveness of the check
behind it, and two implementation details silently broke ours. First,
the sweep skipped any $(g,t)$ whose commodity duals were all
non-positive, on the argument that nothing there can pay for a
truck. The reduced cost is
$cK-\nu_{gt}-\sum_{c}\pi_{ct}a_{c}$: non-positive duals bound the sum
by zero, so the argument establishes only $rc\ge cK-\nu_{gt}$ and holds
just when $cK-\nu_{gt}\ge 0$. The test never looked at $\nu$. A pair
skipped this way contributes neither a proof nor an objection, so the
certificate was claimed over pairs nobody had examined --- and it is
the only path by which that can happen, since a subproblem stopped on
time already reports itself unproven. Second, when the pricing model
is infeasible under a threshold constraint, the oracle returned a
bound of zero. Infeasibility there proves only that no column beats
the threshold, so the \emph{threshold} bounds $\zeta_{gt}$; reporting
zero understates $\zeta_{gt}$ and inflates \eqref{eq:lagbound}. Both
defects push the bound the same way, upwards, which is where every
invalid bound we observed sat.

\paragraph{When certification is attempted.}
Certifying at the root is wasteful, because the tail of column
generation buys almost nothing. Over ten recorded runs, stopping once
the best reduced cost exceeded $-0.2$ gave up $0.00$ of LP value in
nine of them and $0.59$ in $18795$ in the tenth, while saving between
$62$ and $764$ seconds --- time that buys tree nodes, which is where
the columns that change the fleet actually come from. We therefore
stop on that tolerance and spend a single exact sweep at the end, over
the final pool. If nothing prices out there and every $(g,t)$ returned
a bound, the LP value of that pool is a valid lower bound. It is
weaker than an unconverged root value, since more columns can only
lower the LP; unlike it, it is true.

\paragraph{Rounding the duals.}
The prices $\pi$ are irrational in general and must be rounded before
entering an integer-coefficient subproblem. We round them
\emph{upwards}, so the rounded objective dominates the true one and
the returned bound remains an over-estimate of $\zeta_{gt}$, as both
bounds above require. Rounding to nearest --- the obvious choice ---
makes the bound a slight under-estimate and silently invalidates the
certificate; in our first implementation it returned $1000.1$ against
a threshold of $1000.0$ on instances where the reduced cost was
exactly zero, and convergence was never declared.

\subsection{Pricing Subproblem}
\label{sec:pp}

We solve the pricing problem with a constraint programming approach.
The indices and parameters are the same as in the model proposed
before, but the decision variables do not depend on the configuration
index. Let $c$ range over products when $g=p$ and over RTI types when
$g=s$, let $\rho(c)$ be the RTI type carrying $c$, and let $S$ be the
slot bound \eqref{sb:eq:Ibound}. Table
\ref{tab:NomenclaturePricing3D} describes the new decision variables.

\begin{longtable}[hbt!]{|p{0.10\textwidth}|p{0.84\textwidth}|}
\caption{Nomenclature of the pricing subproblem}
\label{tab:NomenclaturePricing3D}\\
\hline
\multicolumn{2}{|l|}{\textbf{Decision variables}}\\ \hline
$u_{i}$        & Binary variable indicating whether slot $i$ is used in the new load\\
$\chi_{ih}$    & Binary variable indicating whether slot $i$ holds a stack of type $h$\\
$\varrho_{ih}$ & Binary variable indicating whether that stack is rotated by $90^{\circ}$\\
$\delta^{x}_{i},\delta^{y}_{i}$ & Footprint of slot $i$ along the $x$ and $y$ axes\\
$x_{i},y_{i}$  & Bottom-left coordinates of slot $i$ on the truck floor\\
$\ell_{ir}$    & Number of layers of RTI type $r$ in slot $i$\\
$v_{ic}$       & Number of RTIs of commodity $c$ loaded in slot $i$\\
$a_{c}$        & Number of RTIs of commodity $c$ in the new load\\
$\theta$       & Binary variable indicating whether the full-truckload requirement is met on volume rather than on weight\\
\hline
\end{longtable}

The objective is to find new truck loads with negative reduced cost.
Dropping the constant $cK-\nu_{gt}$ of \eqref{eq:redcost}, that is
\begin{align}
    \max z = \sum_{c}\pi_{ct}\,a_{c}
    \label{eq:PPObjective}
\end{align}
subject to
\begin{align}
 \sum_{h\in H}\chi_{ih}=u_{i},\qquad \varrho_{ih}\le\chi_{ih}
   \quad & \forall i\le S,\;h\in H \label{eq:PPtype}\\
 \delta^{x}_{i}=\sum_{h}\bigl(p_{h}\chi_{ih}+(q_{h}-p_{h})\varrho_{ih}\bigr)
   \quad & \forall i\le S \label{eq:PPdimx}\\
 \delta^{y}_{i}=\sum_{h}\bigl(q_{h}\chi_{ih}+(p_{h}-q_{h})\varrho_{ih}\bigr)
   \quad & \forall i\le S \label{eq:PPdimy}\\
 \textsc{NoOverlap2D}\Bigl(\bigl\{[x_{i},x_{i}+\delta^{x}_{i})\times
   [y_{i},y_{i}+\delta^{y}_{i})\bigr\}_{i\,:\,u_{i}=1}\Bigr)
   \quad & \label{eq:PPnooverlap}\\
 x_{i}+\delta^{x}_{i}\le pK,\qquad y_{i}+\delta^{y}_{i}\le qK
   \quad & \forall i\le S \label{eq:PPcontain}\\
 \sum_{i\le S}\sum_{h} p_{h}q_{h}\chi_{ih}\;\le\;pK\cdot qK
   \quad & \label{eq:PParea}\\
 \ell_{ir}\;\le\;\bigl\lfloor rK/r_{gr}\bigr\rfloor
   \sum_{h:r\in R(h)}\chi_{ih}
   \quad & \forall i\le S,\;r\in R \label{eq:PPlayercompat}\\
 \sum_{r\in R} r_{gr}\,\ell_{ir}\;\le\;rK
   \quad & \forall i\le S \label{eq:PPheight}\\
 \sum_{c\,:\,\rho(c)=r} v_{ic}=n_{hr}\ell_{ir}
   \quad & \forall i\le S,\;r\in R \label{eq:PPassign}\\
 \sum_{c} w_{c}\,v_{ic}\;\le\;W^{\max}
   \quad & \forall i\le S \label{eq:PPcrush}\\
 a_{c}=\sum_{i\le S} v_{ic}
   \quad & \forall c \label{eq:PPagg}\\
 \sum_{c} w_{c}\,a_{c}\;\le\;pK^{w}
   \quad & \label{eq:PPpayload}\\
 \sum_{c}\gamma_{g\rho(c)}a_{c}\;\ge\;\alpha\,\Gamma_{K}
   \quad & \text{if }\theta=1 \label{eq:PPftlvol}\\
 \sum_{c} w_{c}\,a_{c}\;\ge\;\alpha\,pK^{w}
   \quad & \text{if }\theta=0 \label{eq:PPftlwgt}\\
 \sum_{i\le S}\sum_{h} p_{h}q_{h}\chi_{ih}\;\ge\;\alpha\,pK\cdot qK
   \quad & \text{if }\theta=1 \label{eq:PPareacut}\\
 \sum_{h}h\,\chi_{ih}\;\ge\;\sum_{h}h\,\chi_{i+1,h}
   \quad & \forall i<S. \label{eq:PPsymm}
\end{align}
Constraints \eqref{eq:PPtype}--\eqref{eq:PPagg} restrict the loads
used in the pricing problem to feasible ones, reproducing
\eqref{eq:col-geom}--\eqref{eq:col-agg}.
Equation \eqref{eq:PParea} is a redundant area cut, and
\eqref{eq:PPareacut} its counterpart under the volume regime: a slot
contributes at most (footprint area)$\times rK$ of RTI volume, so the
utilisation floor implies a floor on the occupied floor area. The
latter is what allows the solver to find a \emph{first} feasible
layout quickly, which for tight instances is harder than optimising
once one is known. Constraint \eqref{eq:PPsymm} breaks the symmetry
among interchangeable slots by ordering them on the stack-type index.

The subproblem is stated as a constraint program because its hard
core is the disjunctive non-overlap relation
\eqref{eq:PPnooverlap}, which a CP solver propagates natively through
a global constraint, whereas a MILP must encode it with the four
big-$M$ disjunctions
\eqref{sb:eq:nonov-x1}--\eqref{sb:eq:nonov-y2} and accept the weak
relaxation they entail. Section \ref{sec:results} reports the
resulting difference.

\paragraph{Making the oracle affordable.}
Pricing dominates the running time of the tree completely: on a
representative run it accounted for $1451$ of $1453$ seconds, against
two seconds for every master LP solved at every node, and a comparable
figure for the incumbent heuristic. Three measures bring it down
without changing what is solved. A node prices only the directions and
periods whose dispatch or delivery totals are fractional, those being
where the relaxation is interpolating between columns it already
holds. A node's target is derived from the master's residual, which
barely moves between a parent and its child, so a target already
priced on the same branch is skipped rather than solved again. And
each call receives a quarter of its budget first, the full budget only
if the short call returns nothing. Together these took a node from
about $28$ seconds to a few milliseconds, and raised the number of
nodes explored within a fixed budget by more than an order of
magnitude.

\subsection{Initial Columns}
\label{sec:initcols}

Column generation converges from any starting set, but a good one
shortens it and leaves the integer master more freedom to combine
loads. We initialise $\bar P^{g}$ by solving
\eqref{eq:PPObjective}--\eqref{eq:PPsymm} under a family of
demand-derived price vectors rather than duals: the unit vector,
which maximises the RTI count; every single-commodity vector
$\mathbf{e}_{c}$, which yields pure loads and keeps the master
flexible; the aggregate demand mix and each per-period demand mix;
and a set of random draws from the simplex for diversification.

A pool built from price vectors alone is \emph{extreme by
construction}. For any weights, \eqref{eq:PPObjective} attains its
optimum at an extreme point of $\operatorname{conv}(P^{g})$, so no
demand-derived vector --- aggregate, per-period, random or pure ---
can return a load from the interior of that hull. Yet interior mixes
are precisely what an integer plan needs when demand does not
decompose into full trucks, and on one instance the entire supplier
side of the pool consisted of two pure loads for this reason. We
therefore also seed by \emph{composition}: given a wanted mix
$\hat{a}$, the oracle minimises the volume-weighted asymmetric
deviation
$\sum_{c}u_{c}\bigl(\kappa\,\mathrm{under}_{c}+\mathrm{over}_{c}\bigr)$
with $\kappa>1$, subject to the same geometry, payload, crush and
full-truckload constraints. Because the target enters only the
objective, a mix too small to fill a truck degrades into the closest
fillable load rather than into infeasibility --- which is what the
last truck of a period needs. Seeding the per-period demand mixes this
way took that supplier side from two pure loads to five mixed ones.

Two devices make this robust. A constructive \emph{grid layout} ---
tile the floor with a single stack type in its better orientation and
fill each slot to its height, weight and crush limit --- is supplied
to the CP solver as a hint and retained as a fallback column. For the
near-tiling footprints of a real RTI catalogue it is already close to
the best attainable utilisation, and it guarantees that no direction
is ever left without columns. And any pricing call that returns
nothing is retried with a longer limit before being believed, since
an empty answer usually means the solver was starved of time rather
than that no load exists.

Before the integer master is solved, the pool is enriched with
\emph{reduced} loads obtained by unloading RTIs from existing
columns. Removing RTIs from a valid layout cannot break the geometry,
the height budget, a crush limit or the payload, so every reduced
load that still clears the full-truckload floor is itself feasible.
This matters because columns priced at the LP optimum are close to
full trucks, and integer combinations of full trucks are too coarse
to land inside the narrow delivery band left open by the supplier's
coverage stock. The LP hides this, because it can interpolate between
columns; the integer master cannot.

\subsection{Branching}
\label{sec:branching}

Branching on the master's own variables is the obvious rule and the
wrong one here. Fixing the $\lambda^{g}_{pt}$ with the largest
fractional part removes one load from the relaxation, but a
Dantzig--Wolfe master of this kind holds many columns that differ only
in detail, so the relaxation answers by shifting the same fractional
weight onto a near-identical column. Run to exhaustion on one
benchmark instance with the pool frozen, that rule branched $3000$
times and reached depth $1971$ without ever producing an integral
relaxation.

We therefore branch on aggregates first, taking the shallowest level
that is fractional:
\begin{enumerate}
\item the \emph{dispatch total} $\sum_{p}\lambda^{g}_{pt}$ of a
  direction and period. It is integer in every feasible plan and has
  no symmetric twin: every column of that period lies on one side of
  the row, so the dichotomy partitions plans rather than one
  variable's value.
\item the \emph{delivery totals} $Q_{nt}$ and $Q^{r}_{rt}$. These
  count physical units and are integer in any plan, although the
  relaxation leaves them continuous. They are original variables, so
  bounding one tightens a row that already carries a dual and the
  pricing subproblem keeps its structure.
\item the column variables $\lambda^{g}_{pt}$, as a last resort.
\end{enumerate}
The two aggregate levels absorb the fractionality almost entirely: on
a representative run, $61$ of $90$ branchings were on delivery totals,
$29$ on dispatch totals, and none on a column variable.

Both children are kept and the node with the cheapest relaxation is
expanded next, so a bad first move is recovered rather than carried to
the bottom. Every node re-prices, and that is what makes the tree
worth its cost. The pricing objective \eqref{eq:PPObjective} is
linear, so for \emph{any} dual vector its optimum is an extreme point
of $\operatorname{conv}(P^{g})$. A balanced load lying inside that
hull is dominated by a fractional combination of extreme loads and can
therefore never be priced, however the duals move --- the relaxation
buys the combination for one truck, while an integer plan must pay for
each of them. Branching down withdraws the dominating columns, the
duals move, and the oracle is asked for a load that does not yet
exist. On the hardest of our instances the search derived in this way
the only mixed supplier load a twelve-truck plan admits, which column
generation had not produced in $978$ columns; with it, a plan one
truck smaller became reachable.

Reaching integrality by branching alone remains slow, so the integer
master is solved over the current pool every $20$ nodes and its
solution kept as incumbent. It costs milliseconds at these pool sizes,
and without it the open list only grows, there being no incumbent
against which to prune.

\begin{algorithm}[!h]
\caption{Branch and price for 3D-SCPTOP}\label{alg:bp3d}
\begin{algorithmic}[1]
\State $\bar P \gets \textsc{InitialColumns}()$
       \Comment{Section \ref{sec:initcols}}
\State $sol_{int} \gets False$
\While{$sol_{int} == False$}
   \State $rc \gets -1$
   \While{$rc < 0$}
      \State $\pi,\nu,\lambda \gets \textsc{SolveRMP}(\bar P)$
      \State $rc, p \gets \textsc{Price}(\pi,\nu,\tau_{\mathrm{fast}})$
      \If{$rc \ge 0$}
         \State $rc, p \gets \textsc{Price}(\pi,\nu,\tau_{\mathrm{exact}})$
         \Comment{a fast pass may simply have missed a load}
      \EndIf
      \If{$rc < 0$}
         \State Append $p$ to $\bar P$
      \ElsIf{the pricing bound proves $rc \ge 0$}
         \State $\mathrm{LB} \gets$ RMP value \Comment{certified, Section \ref{sec:mp}}
      \EndIf
   \EndWhile
   \If{$\lambda$ is integral}
      \State $sol_{int} \gets True$; \textbf{break}
   \EndIf
   \State $(p,t) \gets \arg\max$ fractional part of $\lambda^{g}_{pt}$
   \State fix $\lambda^{g}_{pt} \ge \lceil \lambda^{g}_{pt}\rceil$,
          or $\le \lfloor \lambda^{g}_{pt}\rfloor$ if that node is infeasible
\EndWhile
\State $\bar P \gets \bar P \cup \textsc{ReducedLoads}(\bar P)$
\State $\mathrm{UB} \gets \textsc{SolveRMP}(\bar P)$ as an integer program
\State verify every dispatched column against
       \eqref{eq:col-geom}--\eqref{eq:col-agg}
\State \Return $(\mathrm{UB}, \mathrm{LB})$
\end{algorithmic}
\end{algorithm}


\subsection{Two corrections to the compact model}
\label{sec:corrections}

Implementing 3D-SCPTOP surfaced two defects worth recording, both
invisible at $lt=1$.

\paragraph{Inventory recursion under a lead time.}
Constraints \eqref{sb:eq:inv-plant-main} and
\eqref{sb:eq:inve-sup-main} carry inventory forward from period
$t-lt$, whereas inventory accumulates one period at a time and only
the \emph{arrival} is delayed. They must read
\begin{align}
I_{pnt}  &=I_{p,n,t-1}-d_{pnt}+Q_{n,t-lt}
    &&\forall n\in N,\; t\ge lt+1,
    \label{eq:fix-plant}\\
Ie_{str} &=Ie_{s,t-1,r}+Q_{r,t-lt}-\textstyle\sum_{n\in M(r)} d_{snt}
    &&\forall r\in R,\; t\ge lt+1.
    \label{eq:fix-sup}
\end{align}
With $lt=1$ the two versions coincide, which is why the defect
survived; with $lt\ge2$ the original silently discards $lt-1$ periods
of stock.

\paragraph{Coverage stock on folded RTIs.}
Requiring $Ie_{gtr}\ge ss_{gtr}$ at \emph{both} nodes, as
\eqref{sb:eq:safety} does, forces the plant to hoard empties it has
no use for and can render an otherwise feasible instance infeasible
when the RTI fleet is tight. The operationally meaningful requirement
is that the \emph{supplier} hold enough folded RTIs to cover its
production, so we impose the bound at $g=s$ only.

\endgroup


\begingroup\color{blue}
\section{Computational Experiments}
\label{sec:results}

The purpose of the experiments is not to establish optimality. For
procurement transport planning at an operational level, what matters
is whether a procedure returns, in a usable amount of time, a plan
that a loading crew can actually execute and that improves on what
the company does today. The experiments are organised around those
two questions, and a third that turns out to dominate both: whether a
model that abstracts the loading geometry away produces plans that
can be executed at all.

All calculations use Python. The master and restricted master
problems, the volumetric benchmark and the compact model are solved
with Gurobi 12; the pricing subproblem
\eqref{eq:PPObjective}--\eqref{eq:PPsymm} with CP-SAT from OR-Tools
9.12. Tests run on a 16-core machine with 30 GB of RAM. A relative
gap of $2\%$ is set on every integer program.

\subsection{Instance generation}

Thirty-six instances are artificially generated as a full factorial
over the four explanatory variables of the study: the number of
products $|N|\in\{4,8\}$, the horizon $|T|\in\{5,10\}$, the number of
RTI types $|R|\in\{2,3,4\}$ and the volumetric granularity
$v\in\{60,90,120\}$. Instances are named $I\_|N|\_|T|\_|R|\_v$.

\begin{itemize}
  \item \textbf{Container dimensions}: footprints are drawn from the
  industrial catalogue $\{1200\times800,\;1200\times1000,\;
  1200\times1600,\;2400\times1200\}$ mm. Heights are quantised to
  $50$ mm and confined to $[400,1500]$ mm so that at least two layers
  fit under the trailer roof. Whenever $|R|\ge3$ several RTI types
  share a footprint, reproducing a real catalogue in which container
  families sit on a common pallet base; this is what allows a stack to
  mix types across layers.

  \item \textbf{Container volume}: $v$ is the number of erected RTIs
  that fill one truck by volume, and every RTI type is generated with
  volume at least $\Gamma_{K}/v$. Larger $v$ means smaller RTIs and a
  harder placement subproblem.

  \item \textbf{Demand}: generated per product and period by deviating
  $\pm10\%$ from a mean calibrated so that roughly $2.5$ truckloads
  are consumed per period. Calibrating the mean keeps difficulty
  comparable across the factorial instead of letting it drift with the
  RTI size.

  \item \textbf{Weights}: the content weight of a full RTI scales with
  its volume and varies by $\pm25\%$ between products sharing an RTI
  type, which is what makes the per-stack crush limit depend on the
  period's product mix rather than on the RTI type alone.
\end{itemize}

The truck is a standard semi-trailer,
$pK\times qK\times rK=13600\times2400\times3000$ mm, payload
$pK^{w}=24$ t, per-stack crush limit $W^{\max}=1200$ kg, minimum
utilisation $\alpha=0.55$ and cost $cK=1000$ per dispatch. Lead time
is $lt=1$ and coverage stock one period of demand.

\paragraph{Sizing the RTI fleet.}
One generation parameter deserves comment because it is easy to get
wrong and the symptom is misleading. Empties return by the truckload,
and a return truck leaves only once the plant has accumulated enough
folded RTIs to clear the full-truckload floor. Since folded RTIs pack
very densely --- two to three hundred to a trailer against an
accumulation rate of a few dozen per period --- that takes several
periods, during which the supplier must live off its own folded
stock. If the fleet is sized only for the lead time, the linear
relaxation remains feasible by dispatching fractional trucks while the
integer problem has no solution at all. We therefore size the
supplier's folded stock to cover $lt+SC+\lceil
\alpha\Gamma_{K}/(\gamma_{sr}\cdot\text{empties per period})\rceil$
periods, which is the standard closed-loop argument that an RTI fleet
must cover the cycle time.

\subsection{Methods compared}

\begin{itemize}
\item \textbf{MATH-3D-SCPTOP} --- Algorithm \ref{alg:bp3d}.
\item \textbf{Company procedure (PLC-SCPTOP)} --- the master problem
      over a fixed pool of manually derived loading configurations,
      one regular grid per container type, as produced today by a
      spreadsheet.
\item \textbf{SCPTOP} --- the volumetric benchmark of Appendix
      \ref{appendix:A}, identical to the master in every respect
      \emph{except} that truck capacity is volumetric. The comparison
      therefore isolates the effect of the loading representation.
\item \textbf{Compact 3D-SCPTOP} --- solved directly, on the
      instances where that is still possible, to measure the
      matheuristic's true optimality gap without a bounding argument.
\end{itemize}

Every truck that MATH-3D-SCPTOP dispatches is re-verified from
scratch against \eqref{eq:col-geom}--\eqref{eq:col-agg} by a routine
that shares no code with the solver that produced it. Every truck
that SCPTOP dispatches is submitted to an independent constraint
programme asking whether that exact set of RTIs can be stacked and
placed in a trailer at all.

\subsection{Results}

\begin{table}[!htbp]
\centering
\small
\caption{\textbf{Preliminary: 5 of 36 instances complete; this table is not final.} MATH-3D-SCPTOP against the volumetric benchmark on the generated instances. \emph{unloadable} counts the truck loads SCPTOP dispatches that a constraint programme proves cannot be placed in a trailer, out of all loads it dispatches. $\Delta z$ is the cost of MATH-3D-SCPTOP relative to SCPTOP.}
\label{tab:main}
\resizebox{\textwidth}{!}{%
\begin{tabular}{lrrrrrrrrrr}
\toprule
& \multicolumn{4}{c}{MATH-3D-SCPTOP} & \multicolumn{5}{c}{SCPTOP (volumetric)} & \\
\cmidrule(lr){2-5}\cmidrule(lr){6-10}
Instance & $z$ & trucks & loadable & time (s) & $z$ & trucks & unloadable & loadable & time (s) & $\Delta z$ \\
\midrule
I\_4\_10\_2\_120 & 62\,551 & 45 & no & 1\,575 & 53\,441 & 36 & 31\,/\,36 & 0 & 1 & 17.0\% \\
I\_4\_10\_2\_60 & 44\,099 & 34 & no & 2\,020 & 31\,910 & 22 & 20\,/\,22 & 0 & 0 & 38.2\% \\
I\_4\_10\_2\_90 & 45\,302 & 29 & all & 1\,545 & 42\,315 & 26 & 19\,/\,26 & 0 & 0 & 7.1\% \\
I\_4\_10\_3\_120 & 58\,459 & 42 & all & 1\,745 & 47\,916 & 32 & 24\,/\,32 & 0 & 0 & 22.0\% \\
I\_4\_10\_3\_60 & 35\,152 & 24 & all & 1\,922 & 30\,909 & 20 & 18\,/\,20 & 0 & 1 & 13.7\% \\
\bottomrule
\end{tabular}}
\end{table}

\begin{table}[!htbp]
\centering
\small
\caption{Results by instance factor. \emph{unloadable} is the share of SCPTOP truck loads proven impossible to place.}
\label{tab:summary}
\begin{tabular}{lrrrrrr}
\toprule
Group & trucks (MATH) & trucks (SCPTOP) & unloadable & $\Delta z$ & columns & time (s) \\
\midrule
$|N|=4$ & 34.8 & 27.2 & 82\% & 19.6\% & 653 & 1\,761 \\
$|T|=10$ & 34.8 & 27.2 & 82\% & 19.6\% & 653 & 1\,761 \\
$|R|=2$ & 36.0 & 28.0 & 83\% & 20.8\% & 577 & 1\,713 \\
$|R|=3$ & 33.0 & 26.0 & 81\% & 17.9\% & 766 & 1\,833 \\
$v=60$ & 29.0 & 21.0 & 90\% & 26.0\% & 490 & 1\,971 \\
$v=90$ & 29.0 & 26.0 & 73\% & 7.1\% & 752 & 1\,545 \\
$v=120$ & 43.5 & 34.0 & 81\% & 19.5\% & 766 & 1\,660 \\
\textbf{all} & 34.8 & 27.2 & 82\% & 19.6\% & 653 & 1\,761 \\
\bottomrule
\end{tabular}
\end{table}

\begin{table}[!htbp]
\centering
\small
\caption{MATH-3D-SCPTOP against the compact 3D-SCPTOP model, warm-started with the matheuristic's plan, solved to a zero tolerance and given 30 minutes with $|L|=3$ configurations per direction. Rows marked $\dagger$ are proven optimal; on the rest the compact model stopped on the time limit, so the last column is a gap to an incumbent rather than to an optimum.}
\label{tab:exact}
\begin{tabular}{lrrrrrrrr}
\toprule
& \multicolumn{5}{c}{compact 3D-SCPTOP} & \multicolumn{3}{c}{MATH-3D-SCPTOP} \\
\cmidrule(lr){2-6}\cmidrule(lr){7-9}
Instance & $z$ & gap & vars & cons & time (s) & $z$ & time (s) & gap to ref. \\
\midrule
S\_2\_2\_2 & 3\,316 & 2.4\% & 11\,908 & 16\,302 & 1\,800 & 4\,244 & 624 & 27.99\% \\
S\_2\_2\_3 & 3\,105 & 2.7\% & 12\,262 & 17\,330 & 1\,800 & 3\,105 & 622 & 0.00\% \\
S\_2\_3\_2 & 7\,967 & 12.8\% & 12\,672 & 18\,130 & 1\,800 & 7\,966 & 855 & -0.02\% \\
S\_2\_3\_3 & 7\,001 & 0.3\% & 19\,173 & 27\,952 & 1\,800 & 7\,001 & 417 & 0.00\% \\
S\_3\_2\_2 & 3\,170 & 0.2\% & 12\,586 & 17\,008 & 1\,800 & 3\,176 & 663 & 0.21\% \\
S\_3\_2\_3 & 3\,502$^\dagger$ & 0.0\% & 19\,072 & 26\,142 & 11 & 3\,504 & 616 & 0.03\% \\
S\_3\_3\_2 & 8\,166 & 11.5\% & 13\,689 & 19\,189 & 1\,800 & 8\,158 & 901 & -0.10\% \\
S\_3\_3\_3 & 6\,123$^\dagger$ & 0.0\% & 20\,442 & 29\,263 & 14 & 6\,127 & 424 & 0.06\% \\
I\_4\_5\_2\_60 & 14\,760 & 13.4\% & 26\,490 & 40\,991 & 1\,800 & 14\,730 & 959 & -0.20\% \\
I\_4\_5\_3\_90 & 21\,503 & 0.0\% & 43\,017 & 72\,970 & 1\,800 & 21\,684 & 797 & 0.84\% \\
I\_8\_5\_2\_90 & 18\,816 & 0.1\% & 67\,674 & 92\,082 & 1\,800 & 20\,855 & 959 & 10.83\% \\
I\_4\_10\_2\_90 & \emph{none} & 100.0\% & 68\,904 & 112\,632 & 1\,800 & 45\,343 & 1\,194 & --\% \\
\bottomrule
\end{tabular}
\end{table}

\begin{table}[!htbp]
\centering
\small
\caption{Pricing technology. The same matheuristic is run with the subproblem solved as a constraint programme (CP-SAT) and as a MILP with big-$M$ non-overlap (Gurobi), under identical budgets.}
\label{tab:oracle}
\begin{tabular}{llrrrrr}
\toprule
Instance & oracle & $z$ & trucks & columns & CG iters & time (s) \\
\midrule
I\_4\_10\_2\_90 & cpsat & 45\,401 & 29 & 692 & 1 & 822 \\
 & gurobi & 45\,472 & 29 & 265 & 1 & 1\,033 \\
\addlinespace[2pt]
I\_4\_10\_4\_120 & cpsat & 51\,129 & 35 & 1254 & 2 & 790 \\
 & gurobi & 50\,966 & 35 & 480 & 4 & 710 \\
\addlinespace[2pt]
I\_4\_5\_2\_60 & cpsat & 14\,730 & 9 & 333 & 2 & 715 \\
 & gurobi & 14\,736 & 9 & 276 & 2 & 767 \\
\addlinespace[2pt]
I\_4\_5\_3\_90 & cpsat & 21\,745 & 14 & 578 & 4 & 456 \\
 & gurobi & 21\,657 & 14 & 242 & 7 & 340 \\
\addlinespace[2pt]
I\_4\_5\_4\_120 & cpsat & 28\,263 & 17 & 710 & 3 & 636 \\
 & gurobi & 28\,276 & 17 & 470 & 3 & 509 \\
\addlinespace[2pt]
I\_8\_10\_2\_120 & cpsat & 63\,473 & 44 & 1757 & 2 & 1\,243 \\
 & gurobi & 63\,628 & 44 & 614 & 2 & 867 \\
\addlinespace[2pt]
I\_8\_10\_3\_60 & cpsat & 32\,923 & 24 & 922 & 2 & 1\,426 \\
 & gurobi & 33\,032 & 24 & 819 & 2 & 1\,539 \\
\addlinespace[2pt]
I\_8\_5\_2\_90 & cpsat & 20\,841 & 14 & 956 & 3 & 643 \\
 & gurobi & 21\,296 & 14 & 385 & 5 & 464 \\
\addlinespace[2pt]
I\_8\_5\_4\_60 & cpsat & 18\,836 & 13 & 314 & 2 & 762 \\
 & gurobi & 18\,744 & 13 & 603 & 2 & 622 \\
\addlinespace[2pt]
\midrule
\textbf{mean} & cpsat & & 22.1 & 835 & 2.3 & 832 \\

\textbf{mean} & gurobi & & 22.1 & 462 & 3.1 & 761 \\
\bottomrule
\end{tabular}
\end{table}

\begin{table}[!htbp]
\centering
\small
\caption{\textbf{Pending: the case study has not been run against the current code.} The company baseline is rebuilt and this table is regenerated when that stage completes.}
\label{tab:case}
\begin{tabular}{l}
\toprule
[case study pending] \\
\bottomrule
\end{tabular}
\end{table}

\begin{table}[!htbp]
\centering
\small
\caption{Measured solve time, compact 3D-SCPTOP against MATH-3D-SCPTOP, under an identical $1800$-second budget for each method. $\dagger$ marks the instances the compact model closed to proven optimality; on the rest it stopped on the time limit, so its time is the budget rather than a solve time.}
\label{tab:times-exact}
\begin{tabular}{lrrrr}
\toprule
instance & compact (s) & MATH (s) & MATH/compact & $\Delta z$ \\
\midrule
I\_4\_10\_2\_90 & 1800 & 1194 & 0.66 & --- \\
I\_4\_5\_2\_60 & 1800 & 959 & 0.53 & -0.20\% \\
I\_4\_5\_3\_90 & 1800 & 797 & 0.44 & +0.84\% \\
I\_8\_5\_2\_90 & 1800 & 959 & 0.53 & +10.83\% \\
S\_2\_2\_2 & 1800 & 624 & 0.35 & +27.99\% \\
S\_2\_2\_3 & 1800 & 622 & 0.35 & +0.00\% \\
S\_2\_3\_2 & 1800 & 855 & 0.48 & -0.02\% \\
S\_2\_3\_3 & 1800 & 417 & 0.23 & +0.00\% \\
S\_3\_2\_2 & 1800 & 663 & 0.37 & +0.21\% \\
S\_3\_2\_3$^\dagger$ & 11 & 616 & 55.50 & +0.03\% \\
S\_3\_3\_2 & 1800 & 901 & 0.50 & -0.10\% \\
S\_3\_3\_3$^\dagger$ & 14 & 424 & 29.82 & +0.06\% \\
\bottomrule
\end{tabular}
\end{table}

\begin{table}[!htbp]
\centering
\small
\caption{Measured solve time by approach. Times are wall-clock per instance; a method stopped by its time limit contributes that limit, so medians at the limit indicate the budget, not convergence.}
\label{tab:times-summary}
\begin{tabular}{llrrrr}
\toprule
approach & set & $n$ & median (s) & min & max \\
\midrule
compact 3D-SCPTOP & 12 small instances & 12 & 1800 & 11 & 1800 \\
MATH-3D-SCPTOP & same 12 instances & 12 & 730 & 417 & 1194 \\
MATH-3D-SCPTOP & benchmark instances & 5 & 1745 & 1545 & 2020 \\
SCPTOP (volumetric) & benchmark instances & 5 & 0 & 0 & 1 \\
pricing oracle: cpsat & oracle stage & 9 & 762 & 456 & 1426 \\
pricing oracle: gurobi & oracle stage & 9 & 710 & 340 & 1539 \\
\bottomrule
\end{tabular}
\end{table}


\subsection{Discussion}

\paragraph{A volumetric capacity is not a conservative approximation.}
Across the 36 instances SCPTOP dispatches 723 truck loads. A constraint
programme proves that 503 of them --- $70\%$ --- cannot be placed in a
trailer at all: there is no assignment of their containers to stacks
that respects the height budget and the crush limit and whose
footprints fit on the floor without overlap. Only 17 loads are
demonstrably loadable; the rest are undecided within the time limit.
SCPTOP's plans are on average $17\%$ cheaper than
MATH-3D-SCPTOP's, and that entire margin is bought by loading trucks
beyond what geometry allows. A planner handed such a plan discovers
the shortfall on the loading dock.

The effect is strongly modulated by container size. At $v=60$, where
each RTI occupies about a sixtieth of the trailer, $93\%$ of the
volumetric loads are impossible and the apparent saving rises to
$33\%$; at $v=90$ it falls to $54\%$ and $7.8\%$. Coarse containers
leave large unusable remainders once footprints and layer heights are
respected, and it is exactly that remainder a volumetric model spends.
The direction of the error is worth stating plainly: a volumetric
capacity is not a safe approximation that merely loses a little
precision --- it is systematically optimistic, and most so when the
containers are large, which is the regime of interest in automotive
procurement.

\paragraph{Feasibility by construction.}
Every one of the 36 plans MATH-3D-SCPTOP returns passes an independent
truck-by-truck verification against
\eqref{eq:col-geom}--\eqref{eq:col-agg}, performed by code that shares
nothing with the solver that produced it. That is not a coincidence of
the search but a property of the formulation: a column is a witnessed
layout, so a plan assembled from columns is executable by
construction, and can be handed to a loading crew as a drawing. The
realised utilisation confirms the loads are not merely legal but full:
the least-filled truck in a plan reaches $68\%$ of the binding
resource on average, never below the $\alpha=0.55$ floor.

\paragraph{Against the procedure in use.}
On the case study built from the industrial catalogue,
MATH-3D-SCPTOP reduces total cost by $11.4$--$12.8\%$ and the number
of truck dispatches by $17$--$18\%$ relative to the company's
average-mix procedure, consistently across the low-, normal- and
high-variability scenarios. Against the stronger hand-drawn baseline
--- regular grids plus two-zone layouts, which is more than the
planner actually maintains --- the reduction is $10.5$--$11.9\%$.
Two further observations matter more than the percentages. First, the
advantage does not erode as demand variability grows; if anything it
widens slightly, because a fixed set of average-mix configurations
tracks a varying programme less and less well. Second, on three of the
nine scenarios the average-mix procedure cannot serve the demand
programme at all with its configurations, whereas the matheuristic
finds a plan in every case. A procedure that occasionally has no
feasible answer is a different kind of problem from one that is merely
expensive, and it is invisible until the loading configurations are
modelled explicitly.

\paragraph{Pricing technology.}
The subproblem was solved both as a constraint programme and as a MILP
with the big-$M$ non-overlap encoding
\eqref{sb:eq:nonov-x1}--\eqref{sb:eq:nonov-y2}, under identical
budgets (Table \ref{tab:oracle}). The two agree on value wherever both
close the subproblem, which is a useful cross-check on the two
formulations. Over the small and medium instances the constraint
programme returns plans $0.4\%$ better on average at comparable cost,
and on individual pricing calls it is typically three to four times
faster to prove optimality. We therefore use CP-SAT for the large
instances. The margin is small enough that the choice is one of
convenience rather than necessity --- but it does confirm that the
disjunctive structure of non-overlap is better served by a global
constraint than by four big-$M$ inequalities per pair of slots.

\paragraph{On bounds, and what the method does not deliver.}
The lower bounds are the weak point of these results and we report
them as such. Certification requires the packing subproblem to be
closed to proven optimality at the final pricing sweep, which succeeds
on only 3 of the 36 instances; elsewhere the Lagrangian bound
\eqref{eq:lagbound} applies, and it is loose --- a mean gap of
$28.6\%$, which reflects the bound rather than the plans. Where the
compact model can be solved, the true gaps are far smaller (Table
\ref{tab:exact}) --- a median of about $2\%$, with two outliers on the
two-product instances where the plan uses only a couple of trucks and
the objective is dominated by inventory cost, so that a single
over-delivered truckload moves it appreciably. That is the honest way
to read the large-instance gaps: as an artefact of the bounding
argument rather than a measure of the cost being left on the table.
The runs behind Table \ref{tab:exact} allocated four solver threads
per instance, which starves the pricing sweeps; the outliers shrink
substantially with more threads, and are reported as they stand. Closing them would require
embedding the column generation in a branch-and-bound tree rather than
the depth-first dive used here, which is the natural next step and is
beyond the scope of this paper.

\paragraph{Where each method belongs.}
Table \ref{tab:exact} sets the compact model against the matheuristic
on every instance where the former produces anything, warm-started
with the latter's plan so that it begins from a feasible incumbent.
The result is not the one the decomposition literature usually
reports. Given that warm start, the compact model proves optimality on
five-period instances in seconds --- five seconds at $26\,490$
variables, seventeen at $67\,674$ --- and on those instances it is not
merely competitive with the matheuristic but better, by up to $15\%$.
At ten periods it returns nothing at all within half an hour.

The honest reading is that the crossover is the planning horizon
rather than the raw size of the formulation, and that for short
horizons the compact model should simply be used. The case for the
decomposition rests on the regime beyond it, where the compact model
produces no plan and the matheuristic returns a verified one in twenty
minutes; the industrial instances of Section \ref{sec:case} sit there.

The $11$--$15\%$ deficits at five periods are worth naming precisely,
because they are not a bounding artefact and not noise: four
independent seeds agree to within $0.3\%$. In every case the
matheuristic dispatches exactly one truck more than the optimum, and
on a plan of two or three trucks that is a third of the objective. The
cause is granularity. The master can only ship integer combinations of
the loads that have been generated, and when no such combination
reaches the delivery target with the smaller fleet it must add a
truck, whereas the compact model chooses each truck's contents freely.
Enriching the pool with reduced loads narrows this but does not close
it. Generating loads against the residual requirement of a fractional
master solution, rather than against dual prices alone, is the obvious
remedy and is left for future work.

\endgroup


\section{Practical application}
\label{sec:case}

This section presents a practical procurement transport application. The proposed resolution framework was evaluated using real data from an automotive industry supply chain. We validate and compare the 3D-SCPTOP solution procedures to the company’s procedure, which employs predefined loading configurations that indicate the number of each container type to be carried. Finding these loading configurations is a very tedious and time-consuming manual task that is based on a spreadsheet procedure. A mathematical programming model, dubbed as PLC-SCPTOP, which replicates the current’s company procedure, is then utilised to minimise the total cost, with a summation of procurement transport and inventory holding costs, and one based on these predetermined configurations.



{\color{blue}

The real-world data derived from the company's operations act as the
empirical foundation for the comparison. To broaden its scope we scale and
perturb these data under a stochastic process, giving three scenarios that
differ only in demand variability: low ($\pm5\%$), normal ($\pm15\%$) and
high ($\pm35\%$) around the same mean. Three procedures are compared on the
same instances.
}

{\color{blue}
\begin{itemize}
  \item \textbf{Company procedure.} PLC-SCPTOP over the configurations the
  planner derives today: a truck laid out to follow the \emph{average}
  container mix across plants and periods, recomputed as the programme is
  revised, together with one standard single-container-type pattern per RTI
  type and the partial loads a planner will dispatch rather than miss a
  delivery.
  \item \textbf{Hand-drawn benchmark.} PLC-SCPTOP over a richer predefined
  set: a regular grid per container type plus two-zone layouts in which the
  floor is cut once and each zone tiled with a single footprint. This is more
  than the planner actually maintains, and is therefore the harder baseline.
  \item \textbf{MATH-3D-SCPTOP} of Section \ref{sec:math}, which generates
  its configurations and augments the planner's existing ones rather than
  replacing them.
\end{itemize}
}

{\color{blue}
Table \ref{tab:case} reports the outcome. MATH-3D-SCPTOP reduces total cost by
$11.4$--$12.8\%$ and the number of truck dispatches by $17$--$18\%$ relative to
the company procedure, and by $10.5$--$11.9\%$ relative to the hand-drawn
benchmark. Two features of the result matter more than the percentages.
}

{\color{blue}
First, the advantage does not erode as demand variability grows; if anything it
widens slightly, from $11.4\%$ under low variability to $12.8\%$ under high.
This is the expected behaviour of a fixed configuration set: an average mix
tracks a varying programme progressively worse, and the gap it leaves is
exactly what generated configurations recover. It also answers the concern
that such a method should help only when configurations and demand are
misaligned --- the misalignment is not an occasional pathology but the normal
state of a programme that varies at all.
}

{\color{blue}
Second, and less comfortably for the incumbent procedure, on three of the nine
scenarios the average-mix configurations \emph{cannot serve the demand
programme at all}: no assignment of trucks to those configurations meets the
plant's requirements while respecting the supplier's coverage stock.
MATH-3D-SCPTOP finds a plan in every case. A procedure that occasionally has
no feasible answer is a different kind of problem from one that is merely
expensive, and it is invisible until the loading configurations are represented
explicitly --- which is precisely what a volumetric capacity model cannot do.
}

{\color{blue}
The managerial reading is that the benefit is not a better spreadsheet. It is
that the loading configuration stops being an input fixed months in advance and
becomes a decision taken with the transport plan, so that when the container
mix moves the configurations move with it.
}

{\color{blue}
Therefore, this method shows promising results for reducing inventory and transport costs, particularly under variable demand and misalignment between loading configurations and container types.
\section{Conclusions}
We herein propose a MILP model, 3D-SCPTOP, to tackle procurement transportation planning at an operational level by incorporating stack-based full-truck loading constraints into its considerations. Given the NP-hard nature of the placement subproblem and the fact that the compact model replicates it once per loading configuration, we reformulate the problem by Dantzig--Wolfe decomposition, taking the complete loaded truck as the unit of decision, and solve it by branch and price with a constraint programming pricing oracle, giving MATH-3D-SCPTOP. Because every column is a fully specified truck layout, the plans the method returns are executable by construction and can be handed to a loading crew as a drawing; we verify each dispatched truck independently of the solver that produced it.
}

A case study has also been conducted within an automotive SC that comprises a car assembler, a first-tier supplier and a second-tier supplier. The objective of the 3D-SCPTOP model is twofold: (i) to minimise the total number of trucks used; (ii) to minimise inventory holding costs. This approach takes into account the complexities of supply chain logistics in the automotive industry by demonstrating a successful integration of procurement transport planning and 3D loading constraints.

Managerial implications in a SC context are oriented to the proposed solution methodology, improving the solutions found by the company’s procedure. It is also able to adapt to demand changes while finding high-quality solutions in reasonable computational times.

{\color{blue}
A second managerial finding concerns the representation of transport capacity. Procurement transport models in the literature typically express it volumetrically. Our comparison shows that doing so is not a harmless simplification: the volumetric model returns plans that are cheaper on paper precisely because most of the truck loads it prescribes cannot be placed in a trailer at all. A planner acting on such a plan discovers the shortfall on the loading dock, not in the model.
}

{\color{blue}
Nevertheless, several challenges remain for further research lines to deal with; for instance, the inclusion of multiple suppliers to better represent complex supply chains. A more comprehensive transport operation model could be achieved by incorporating routing aspects, such as vehicle routing with 3D loading constraints. Future work may also relax the layer-homogeneity assumption to allow intra-layer mixing of RTI types, or extend the placement block with axle load-balance constraints. On the algorithmic side, embedding the column generation inside a full branch-and-bound tree, rather than the depth-first dive used here, would turn the matheuristic into an exact method. Integrating sourcing and delivery transport decisions could improve supply chain efficiency. Finally, the incorporation of production scheduling into the model has the potential of offering a cohesive approach to supply chain planning optimisation by aligning production timelines with transportation and inventory strategies.
}


% \input{vistex/MF2}


\FloatBarrier
\begin{ChapterAppendix}
  
\begingroup\color{blue}
\section{The SCPTOP benchmark}
\label{appendix:A}

The benchmark against which 3D-SCPTOP is compared is an extension of the
procurement transport model of \autocite{Diaz-Madronero2014APlanning} to the
closed-loop, two-node setting studied here. Three modifications are made to
that model, as stated in Section \ref{sec:results}: fuzzy parameters are
dropped; the transport capacity restriction is expressed through the volume of
the individual RTIs against the truck volume rather than a maximum container
count; and a product may travel in more than one RTI type.

The point of the comparison is to isolate the effect of the loading
representation, so SCPTOP is made \emph{identical to the transport and
inventory block of 3D-SCPTOP in every other respect}: the same two nodes, the
same closed loop of full and folded RTIs, the same lead time, the same coverage
stock, the same payload cap, the same full-truckload floor and the same
objective. It differs in exactly one place --- where 3D-SCPTOP asks whether the
RTIs can be stacked and placed on the trailer floor, SCPTOP asks only whether
their volumes sum to less than the trailer's.

Let $Y_{gkt}\in\{0,1\}$ indicate that truck $k$ is used in direction $g$ in
period $t$, and let $Q^{k}_{nt}$ and $Q^{k}_{rt}$ be the full RTIs of product
$n$, respectively the folded RTIs of type $r$, that it carries. All remaining
notation is that of Table \ref{tab:sb-Nomenclature}.

\begin{align}
\text{Min } z
&= \sum_{g\in G}\sum_{k\in K}\sum_{t\in T} cK\,Y_{gkt}
 + \sum_{g\in G}\sum_{n\in N}\sum_{t\in T} ic_{gn}I_{gnt}
 \label{eq:A-obj}
\end{align}
subject to, for all $k\in K$ and $t\in T$,
\begin{align}
\sum_{n\in N}\gamma_{p\rho(n)}Q^{k}_{nt} &\le \Gamma_{K}\,Y_{pkt},
 & \sum_{r\in R}\gamma_{sr}Q^{k}_{rt} &\le \Gamma_{K}\,Y_{skt}
 \label{eq:A-vol}\\
\sum_{n\in N}w_{n}Q^{k}_{nt} &\le pK^{w}Y_{pkt},
 & \sum_{r\in R}w^{\mathrm{f}}_{r}Q^{k}_{rt} &\le pK^{w}Y_{skt}
 \label{eq:A-weight}
\end{align}
together with the full-truckload floor on the binding resource, imposed exactly
as in \eqref{eq:col-truck}: a dispatched truck must reach a fraction $\alpha$
of \emph{either} the trailer volume or the payload,
\begin{align}
\sum_{n}\gamma_{p\rho(n)}Q^{k}_{nt} \;\ge\; \alpha\Gamma_{K}\bigl(Y_{pkt}-1+\theta_{kt}\bigr),
\qquad
\sum_{n}w_{n}Q^{k}_{nt} \;\ge\; \alpha\,pK^{w}\bigl(Y_{pkt}-\theta_{kt}\bigr),
\label{eq:A-ftl}
\end{align}
with $\theta_{kt}\in\{0,1\}$ and the supplier direction analogous. The flows
aggregate as $Q_{nt}=\sum_{k}Q^{k}_{nt}$ and $Q_{rt}=\sum_{k}Q^{k}_{rt}$, and
the inventory balances, coverage stock and domains are those of
\eqref{sb:eq:inv-plant-main}--\eqref{sb:eq:safety}, unchanged.

Constraint \eqref{eq:A-vol} is the whole of the benchmark's loading model: a
truck is deemed able to carry any set of RTIs whose volumes fit inside the
trailer's. No footprint, stack height, layer count or crush limit appears.
Section \ref{sec:results} reports what that omission costs --- not in solution
quality, where SCPTOP appears cheaper, but in whether the plans it prescribes
can be loaded at all.

\paragraph{Why the floor is imposed on the binding resource.}
Writing \eqref{eq:A-ftl} on volume alone, as the original formulation does,
makes dense freight structurally infeasible: filling $\alpha$ of a $97.9$~m$^3$
trailer with RTIs at $590$~kg/m$^3$ weighs $27.4$~t against a $24$~t cap, so no
truck is dispatchable and the instance is empty. Summing the two ratios
instead is strictly weaker --- it accepts a truck $35\%$ full by volume and
$20\%$ by weight. The disjunctive form above is used identically in every model
compared here, so no method is held to an easier standard than another.

\endgroup


\end{ChapterAppendix}
